% !TeX root = ../../Book1.tex
\OptionalChapter{Sobolev 容量与精细性质}{b1:ch:26}
\BookChapterNumber{b1:ch:26}
\begin{chapterabstract}
本章研究 Sobolev 等价类能够在多小的例外集之外确定点值。以邻域上的最小能量定义容量后，证明它的次可加性，构造拟连续代表，并建立拟处处唯一性及范数收敛的逐点子列结论。容量与 Hausdorff 测度的比较进一步给出异常集的维数限制。最后介绍细拓扑及 Wiener 判据。基础容量性质在正文中证明；涉及非线性势论的深层等价放在选读部分，只给出证明概要。
\end{chapterabstract}

\begin{learninggoals}
\begin{enumerate}
\item\label{b1:goal:26-capacity}从开邻域上的几乎处处约束定义非齐次容量，证明单调性、可列次可加性，并辨别零测与零容量。
\item\label{b1:goal:26-representative}用快速光滑逼近构造拟连续代表，证明超水平容量估计及拟处处唯一性，不用尚未建立的点值概念反过来定义容量。
\item\label{b1:goal:26-dimension}通过小球覆盖、弱紧性与平均振荡估计，把容量零集与 Hausdorff 维数联系起来，单独处理临界指数和对数尺度。
\item\label{b1:goal:26-fine}区分全空间非齐次容量与相对纯梯度容量，理解薄性和细拓扑的定义，以及 Wiener 边界判据还需要哪些势论工具。
\item\label{b1:goal:26-boundaries}辨认拟连续、通常连续与 Lebesgue 精确代表之间已证的关系，避免把几乎处处结论未经证明提升为拟处处结论。
\end{enumerate}
\end{learninggoals}

前几章在积分范数中讨论弱导数、逼近和紧性，函数始终先作为几乎处处等价类出现。第24章在 \(p>d\) 时已经得到唯一连续代表，但当 \(p\le d\) 时，普通连续性不再普遍成立。另一方面，完全放弃点值又会掩盖零边界条件、薄集合和异常集的几何。容量提供了介于这两种处理之间的方式：不要求每个点都被确定，却要求例外集比任意 Lebesgue 零测集更小。

\ProseParagraphBreak

本章属于选读。主线固定 \(1<p<\infty\)，因为次可加性和若干极限步骤使用自反性；端点 \(p=1\) 需要另一套细节，本章不作讨论。前三节主要依赖第9章的 Mazur 定理与凸泛函弱下半连续性、第13章的 \(5r\) 覆盖、第15章的 Hausdorff 测度、第16章和第21章的一维绝对连续理论与 ACL 代表，以及第24章的 Poincaré 和 Morrey 不等式。弱导数、光滑稠密性与 Sobolev 弱紧性仍是贯穿其中的基础。

\ProseParagraphBreak

最后一节把这些结果放入非线性势论的图景。Wiener 判据、Cartan 性质及超调和函数刻画均作为外部输入，并注明出处；它们不由前三节推出。

\ProseParagraphBreak

Kinnunen 与 Martio 的《The Sobolev capacity on metric spaces》\cite{KinnunenMartio1996Capacity}适合对照容量性质和快速逼近的组织；其中的度量空间范数与这里的经典能量仅等价，正文会指出具体差别。系统的欧氏精细性质可继续参阅 Ziemer 的《Weakly Differentiable Functions》\cite{Ziemer1989Weakly}；容量与非线性椭圆方程之间的联系，可伴读 Heinonen、Kilpeläinen 与 Martio 的《Nonlinear Potential Theory of Degenerate Elliptic Equations》\cite{Heinonen2006Nonlinear}。这些专著的范围超过本章，不要求读者在开始阅读前掌握全部势论。

% !TeX root = ../../Book1.tex
\section{\texorpdfstring{\(p\)}{p}-容量}
\label{b1:sec:2601}

一个集合的 Lebesgue 测度为零，意味着在它上面更改函数值不影响任何 \(L^p\) 等价类。但如果要求函数在这个集合的整个邻域中保持某种大小，更改就不再是免费的：过渡到邻域外部需要弱梯度，而梯度也计入 Sobolev 范数。容量用完成这种过渡所需的最小能量度量集合。

本章固定 \(d\ge1\)、\(1<p<\infty\)，容许函数先取实值。这个范围使我们可以使用已经证明的自反性；它不包括 \(p=1\) 的另行处理，也不把 \(p=\infty\) 纳入同一能量定义。记
\begin{equation}\label{b1:eq:sobolev-capacity-energy}
 E_p(u)=\int_{\mathbb R^d}
       \bigl(|u(x)|^p+|\nabla u(x)|^p\bigr)\,\mathrm dx.
\end{equation}
梯度取欧氏长度，故 \(E_p(u)^{1/p}\) 是一个旋转不变的范数，与第21章逐分量 \(p\) 次和定义的标准范数等价；当 \(p\ne2\) 时，二者通常不相等。保留零阶项，使这里的容量成为非齐次全空间容量。\secref{b1:sec:2604}的相对纯梯度容量将另行定义。

Kinnunen 与 Martio 的《The Sobolev capacity on metric spaces》\cite[Sections 1 and 3, p. 368]{KinnunenMartio1996Capacity}提供开邻域容量及其基本性质的组织方式。其主体使用度量空间中的另一种 Sobolev 范数，并明确区分它与经典弱梯度能量；本节直接使用\eqref{b1:eq:sobolev-capacity-energy}，以本书已经建立的截断、弱紧性和弱下半连续性完成证明。
% 正式论文368页(1.3)--(1.5)及371--375页3.1--3.7、4.1已由协作撰写者实读。
% 本节经典欧氏格能量路线独立展开，不把Hajlasz型梯度当作普通弱梯度。

\subsection{定义}
\label{b1:subsec:260101}

\begin{definition}\label{b1:def:sobolev-capacity}
对开集 \(G\subseteq\mathbb R^d\)，定义
\[
 C_p(G)=\inf\left\{E_p(u):
 u\in W^{1,p}(\mathbb R^d),\
 u\ge1\text{ 几乎处处于 }G\right\}.
\]
若容许类为空，规定下确界为 \(+\infty\)。对任意集合 \(A\subseteq\mathbb R^d\)，定义其 \term[sobolev-rongliang]{Sobolev 容量}[Sobolev capacity]，亦称 \(p\)-容量，为
\begin{equation}\label{b1:eq:sobolev-outer-capacity}
 C_p(A)=\inf\{C_p(G): A\subseteq G,\ G\text{ 开}\}.
\end{equation}
\end{definition}
\symidx{\(C_p(A)\)：非齐次 Sobolev \(p\)-容量}

定义分为两步。第一步只要求开集上的几乎处处不等式，因此容许性不依赖 Sobolev 代表的选择。第二步以真正包含 \(A\) 的开集作外逼近，所以即使 \(A\) 本身不可测，也仍有明确的容量。开集的单调性直接来自容许类的包含关系；因此当 \(A\) 本身开时，第二步所得数值与第一步相同。

不能把上述条件改成“选某个代表使其在 \(A\) 上至少为一”。若 \(A\) 是零测集，零函数类可以任意在 \(A\) 上改成一；这种定义会令全部零测集都没有代价，失去我们要研究的区别。也不能一开始就要求“拟处处至少为一”，因为拟处处正要依赖容量才能定义。

截断使容许函数具有统一的数值范围。以下格运算的能量公式也将保证我们能处理多个集合。

\begin{lemma}\label{b1:lem:sobolev-lattice-energy}
若 \(u,v\in W^{1,p}(\mathbb R^d)\) 为实值，则 \(u\vee v=\max(u,v)\)、\(u\wedge v=\min(u,v)\) 都属于该空间，且
\begin{equation}\label{b1:eq:sobolev-lattice-energy}
 E_p(u\vee v)+E_p(u\wedge v)=E_p(u)+E_p(v).
\end{equation}
此外，\(T(u)=\min(1,\max(0,u))\) 满足 \(E_p(T(u))\le E_p(u)\)。若 \(u\) 在开集 \(G\) 上容许，则 \(T(u)\) 仍容许，且 \(0\le T(u)\le1\)，在 \(G\) 上几乎处处等于一。
\end{lemma}
\begin{proof}
由第21章的截断公式，\(u\vee v=v+(u-v)_+\) 属于 Sobolev 空间。在 \(\{u>v\}\) 上，它的弱梯度为 \(\nabla u\)；在 \(\{u<v\}\) 上为 \(\nabla v\)。在相等集合上，\cref{b1:prop:sobolev-gradient-on-level-sets}给出 \(\nabla(u-v)=0\) 几乎处处，两个选择也一致。最小值的公式相反。

于是对几乎所有点，最大值与最小值只是重新排列了两个函数值和两个梯度向量。分别取绝对值或欧氏长度的 \(p\) 次幂后相加，再积分，就得\eqref{b1:eq:sobolev-lattice-energy}。截断的值不增大绝对值，梯度只在 \(0<u<1\) 上保留原梯度，故能量不增。最后的容许性由截断在 \([1,\infty)\) 上恒等于一立即得到。
\end{proof}

\subsection{单调性}
\label{b1:subsec:260102}

\begin{proposition}\label{b1:prop:sobolev-capacity-basic}
容量具有以下性质：
\begin{enumerate}
\item \(C_p(\varnothing)=0\)，且 \(A\subseteq B\) 蕴含 \(C_p(A)\le C_p(B)\)。
\item 容量满足\eqref{b1:eq:sobolev-outer-capacity}的开外正则性，并且
\[
 m_d^*(A)\le C_p(A)
 \quad\text{对任意 }A\subseteq\mathbb R^d.
\]
\item 每个有界集具有有限容量。平移或正交变换不改变容量。
\end{enumerate}
\end{proposition}
\begin{proof}
零函数在空集上容许，所以空集容量为零。开集之间的单调性已由容许类包含得到；再对开邻域取下确界，便得任意集合的单调性。

若 \(A\subseteq G\) 且 \(u\) 在开集 \(G\) 上容许，则
\[
 m_d^*(A)\le m_d(G)\le\int_G|u|^p
 \le E_p(u).
\]
先对容许函数、再对开邻域取下确界，即得外测度下界。开外正则性本身就是第二步定义。

有界集包含在一个球内，取在该球附近等于一的光滑紧支集函数，其有限能量给出容量上界。平移、正交变换保持 Lebesgue 测度和欧氏梯度长度，把一个集合的容许函数与开邻域双向变换，就得到容量不变。
\end{proof}

零阶项使一个简单结论变得明显：无限 Lebesgue 测度的集合具有无限容量。这与只看梯度的全空间齐次问题不同。另一方面，面积为零的薄片仍可能具有正容量，稍后会给出直接证明。

若把 \(E_p\) 换成标准 Sobolev 范数的 \(p\) 次幂，由两种能量的正常数比较，所得容量也双向受常数控制。因此零容量集相同；\secref{b1:sec:2602}按任意小容量开集定义的拟连续性也相同。但对于一个具体集合，容量的数值可能改变，不能在精确常数的计算中随意更换能量。

\subsection{次可加性}
\label{b1:subsec:260103}

本节所说的可列次可加性 (countable subadditivity) 是并集容量不超过各项容量之和；下文简称次可加性。它不同于互不相交并上的可列可加等式。把两个容许函数相加虽然能覆盖两个集合，却会在能量中产生不必要的幂次常数。取最大值并使用格能量公式，可以保留系数一；处理一列集合时，再用弱紧性通过极限。

\begin{theorem}\label{b1:thm:sobolev-capacity-subadditive}
对任意集合列 \((A_j)_{j\ge1}\)，
\begin{equation}\label{b1:eq:capacity-countable-subadditivity}
 C_p\left(\bigcup_{j=1}^\infty A_j\right)
 \le\sum_{j=1}^\infty C_p(A_j).
\end{equation}
\end{theorem}
\begin{proof}
若右端无限，结论显然。以下设其有限，固定 \(\varepsilon>0\)。逐项选择开邻域 \(G_j\supseteq A_j\) 及截断后的容许函数 \(u_j\)，使
\[
 0\le u_j\le1,\qquad
 E_p(u_j)\le C_p(A_j)+\varepsilon2^{-j}.
\]
这里可在选择开集与选择容许函数时分别分配半份误差。令
\[
 v_N=\max(u_1,\ldots,u_N),\qquad v=\sup_j u_j.
\]
反复应用格能量公式可得
\[
 E_p(v_N)\le\sum_{j=1}^N E_p(u_j)
 \le\sum_jC_p(A_j)+\varepsilon.
\]
同时 \(0\le v_N\uparrow v\)，且 \(v^p\le\sum_j|u_j|^p\)，右端可积。控制收敛因此给出 \(v_N\to v\) 于 \(L^p\)。

有界能量等价于 Sobolev 范数有界。由于 \(1<p<\infty\)，可抽取 Sobolev 弱收敛子列；其零阶弱极限由已经取得的 \(L^p\) 强极限确定，必为 \(v\)。所以 \(v\in W^{1,p}\)。能量是凸且范数连续的泛函，故由第9章的凸泛函弱下半连续性，
\[
 E_p(v)\le\liminf_N E_p(v_N)
 \le\sum_jC_p(A_j)+\varepsilon.
\]
每个 \(u_j\) 在 \(G_j\) 上几乎处处等于一。合并可数个例外零集后，\(v\) 在开集 \(\bigcup_jG_j\) 上几乎处处等于一，因此是这个开并的容许函数。该开并包含 \(\bigcup_jA_j\)，故得到\eqref{b1:eq:capacity-countable-subadditivity}，最后让 \(\varepsilon\downarrow0\)。
\end{proof}

至此，\(C_p\) 具有外测度的三个公理。不过它一般不是通常意义的可加测度；不能因其定义在全部集合上，就认定所有 Borel 集自动满足 Carathéodory 可测性。第\ref{b1:ex:26-one-dimensional-exact-capacity}题将具体算出可加性的失败。我们在后续只使用已经证明的单调性、次可加性和开外正则性。

下一条估计把能量与点值联系起来，但暂时只对连续函数陈述。

\begin{lemma}\label{b1:lem:continuous-capacity-level-set}
设 \(v\in W^{1,p}(\mathbb R^d)\cap C(\mathbb R^d)\)，\(t>0\)。则
\[
 C_p(\{|v|>t\})\le t^{-p}E_p(v).
\]
\end{lemma}
\begin{proof}
连续性使超水平集开；函数 \(\min(1,|v|/t)\) 在其上几乎处处等于一，且能量不超过 \(t^{-p}E_p(v)\)，因而是所需容许函数。
\end{proof}

这不是对任意代表都成立的 Chebyshev 不等式。若只知道 \(v\) 是一个可测代表，超水平集未必开，而且在零测集上任意改变它会改变该集合的容量。下一节先构造拟连续代表，再为那一类代表证明对应估计。

\subsection{容量零集}
\label{b1:subsec:260104}

若 \(C_p(N)=0\)，则称 \(N\) 为容量零集。它的任意子集仍为容量零集，可数个容量零集的并也为容量零集；这些集合又都是 Lebesgue 零测集。为了了解这个例外集尺度到底有多细，先考察小球与单点。

\begin{proposition}\label{b1:prop:sobolev-capacity-balls}
存在只依赖 \(d,p\) 的常数，使所有 \(x\in\mathbb R^d\)、\(r>0\) 满足
\[
 C_p(B(x,r))\le C_{d,p}(r^d+r^{d-p}).
\]
\end{proposition}
\begin{proof}
固定 \(0\le\eta\le1\) 的光滑函数，等于一于 \(B(0,1)\) 的一个邻域，且支集包含于 \(B(0,2)\)。函数 \(\eta((\,\cdot-x)/r)\) 在 \(B(x,r)\) 上容许，变量代换给出其零阶能量为 \(r^d\|\eta\|_p^p\)，梯度能量为 \(r^{d-p}\|\nabla\eta\|_p^p\)。两者相加即得。
\end{proof}

\begin{proposition}\label{b1:prop:sobolev-point-capacity}
若 \(1<p\le d\)，每个单点以及每个可数集的 \(p\)-容量均为零。
若 \(p>d\)，存在 \(c_{d,p}>0\)，使每个非空集合 \(A\) 都满足 \(C_p(A)\ge c_{d,p}\)。
\end{proposition}
\begin{proof}
当 \(p<d\) 时，小球上界在 \(r\downarrow0\) 时趋零，单点结论由单调性得到，再用次可加性处理可数并。

当 \(p=d>1\) 时，普通缩放的梯度能量不趋零，需要较慢的对数过渡。固定 \(R>0\)，取 \(0<\varepsilon<R\)，定义径向函数
\[
 v_\varepsilon(x)=
 \begin{cases}
 1,&|x|\le\varepsilon,\\
 \displaystyle\frac{\log(R/|x|)}{\log(R/\varepsilon)},
       &\varepsilon<|x|<R,\\
 0,&|x|\ge R.
 \end{cases}
\]
对每个固定的 \(\varepsilon\)，这是紧支集 Lipschitz 函数，故属于 \(W^{1,d}\)。记 \(\sigma_{d-1}=d\omega_d\) 为单位球面的面积，径向积分给出
\[
 \int|\nabla v_\varepsilon|^d
 =\frac{\sigma_{d-1}}{\log(R/\varepsilon)^d}
            \int_\varepsilon^R\frac{\mathrm dr}{r}
 =\sigma_{d-1}\log(R/\varepsilon)^{1-d}\longrightarrow0.
\]
同时 \(0\le v_\varepsilon\le\mathbf1_{B(0,R)}\)，并在每个非零点趋于零，控制收敛给出其零阶能量趋零。它在原点的开邻域上容许，因此 \(C_d(\{0\})=0\)。平移及次可加性完成本情形。

最后设 \(p>d\)。由 Morrey 估计与范数等价性，每个 Sobolev 函数都有连续代表，并满足
\[
 |u^*(x)|\le C_{d,p}E_p(u)^{1/p}
 \quad(x\in\mathbb R^d).
\]
若 \(u\) 在含 \(x\) 的开集上几乎处处至少为一，连续性说明 \(u^*(x)\ge1\)：否则某个小球上都小于一，与几乎处处条件矛盾。所以每个这样的容许函数都有统一的正能量下界。两次取下确界得到单点正下界，再用单调性得到全部非空集合的下界。
\end{proof}

在 \(d\ge2\) 时，一个平面片给出更直观的零测而正容量集合。令
\[
 S=[0,1]^{d-1}\times\{0\}.
\]
若开集 \(G\) 包含 \(S\)，紧性保证存在 \(\delta>0\)，使 \([0,1]^{d-1}\times(-\delta,\delta)\subset G\)。对在 \(G\) 上容许的 \(u\)，ACL 性质与 Fubini 定理表明，对几乎所有 \(x'\in(0,1)^{d-1}\)，竖线上的一维绝对连续代表在高度零处至少为一。一维点值估计给出
\[
 1\le C_p\int_{\mathbb R}
          \bigl(|u(x',t)|^p+|\partial_du(x',t)|^p\bigr)\,\mathrm dt.
\]
为看出该估计的常数不依赖 \(\delta\)，可任取 \(s\in(0,1)\)，由
\[
 |u(x',0)|\le |u(x',s)|+\int_0^1|\partial_du(x',t)|\,\mathrm dt
\]
对 \(s\) 平均，再用 Hölder 不等式。随后对 \(x'\) 积分，得到 \(1\le C_p E_p(u)\)。因此 \(C_p(S)>0\)，尽管 \(m_d(S)=0\)。邻域厚度可以任意小，保持片上值所需的能量却不能消失。

这些例子揭示了容量与代表理论之间的联系。超过维数的 Sobolev 指数已经使点值连续受控，于是一个点也不能被视为可忽略的容量例外。指数不超过维数时，单点可以忽略，但不是所有 Lebesgue 零测集都可以忽略；下一节将在这个更严格的例外集尺度下讨论代表的唯一性。

% !TeX root = ../../Book1.tex
\section{拟处处性质}
\label{b1:sec:2602}

弱导数只由几乎处处等价类决定，而点值需要指定代表。容量提供了比 Lebesgue 零测集更细的例外集尺度：我们希望在这个尺度下选出合适的代表，并说明哪些逐点结论不再依赖选择。本节固定 \(1<p<\infty\)，先讨论实值 \(W^{1,p}(\mathbb R^d)\)，容量及能量 \(E_p\) 均沿用\cref{b1:def:sobolev-capacity}。特别地，\(E_p^{1/p}\) 与通常 Sobolev 范数等价，梯度项用欧氏模。

\ProseParagraphBreak

Kinnunen 与 Martio 的《The Sobolev capacity on metric spaces》\cite[Section 3, 3.6. Theorem and 3.7. Corollary, p. 368]{KinnunenMartio1996Capacity}用可和的坏集容量，从连续函数逼近构造拟连续代表。这条路线也适用于这里已经建立的经典容量。该文研究的度量空间范数与本书能量不同，并明确说明，在欧氏空间中两种范数及相应容量等价，而非数值相同。

\ProseParagraphBreak

“拟处处”与“拟连续”的中文名称可对照王巍翰、赵唯的《Finsler 度量测度空间上的容量》\cite[定义4.4]{WangZhao2023Capacity}。那里允许删去一般的小容量集；借助本章的开外正则性，可以把它包含在容量仍任意小的开集中。本书在定义时直接要求坏集开放，以便随后使用闭补集上的相对连续性；名称对照不意味着把其 Finsler 空间结论未经证明移到这里。

% 实读正式期刊公开全文 https://www.acadsci.fi/mathematica/Vol21/kinnunen.pdf：
% 第368页(1.3)--(1.5)区分Hajlasz型范数与经典欧氏能量；
% 第371--374页3.1--3.7，重点3.6的快速逼近完整证明及3.7；
% 第375页4.1的测度下界。第374页另核看实际页面。
% 本节不照搬该文容量数值；拟连续超水平估计及唯一性按经典开集容量独立证明。
% 前置：26.1格能量、次可加性、连续函数超水平估计、m*<=C；
% 22.1全空间光滑稠密。未调用一般精细拓扑或容量版Lebesgue点定理。

\subsection{拟处处}
\label{b1:subsec:260201}

\begin{definition}\label{b1:def:capacity-quasi-everywhere}
设 \(E\subseteq\mathbb R^d\)。若某性质在 \(E\) 上除去一个 \(C_p\)-容量为零的集合后成立，则称它在 \(E\) 上\term[nichuchu]{拟处处}[quasi-everywhere]成立，记作 \(p\)-q.e.；固定 \(p\) 后也简写为 q.e.
\end{definition}
\symidx{\(p\text{-q.e.}\)：除去 \(p\)-容量零集后成立}

由\cref{b1:prop:sobolev-capacity-basic}中的 \(m_d^*(N)\leq C_p(N)\)，容量零集一定是 Lebesgue 零测集，所以拟处处成立蕴含几乎处处成立。这是容量估计的后果，不是两种例外集的定义相同。反向蕴含不能由此推出。

\ProseParagraphBreak

可数个拟处处结论可以放在同一个例外集外使用。事实上，若第 \(j\) 条结论的例外集为 \(N_j\)，则次可加性给出
\[
 C_p\left(\bigcup_{j=1}^{\infty}N_j\right)
 \leq\sum_{j=1}^{\infty}C_p(N_j)=0.
\]
同样，如果一列集合 \(A_j\) 满足 \(\sum_j C_p(A_j)<\infty\)，则
\begin{equation}\label{b1:eq:capacity-summable-exceptions}
 C_p\left(\bigcap_{N=1}^{\infty}\bigcup_{j\geq N}A_j\right)
 \leq\inf_N\sum_{j\geq N}C_p(A_j)=0.
\end{equation}
这只用了单调性与次可加性，没有把容量当成可数可加测度。其含义是：拟处处的点只会落入有限多个 \(A_j\)。

\ProseParagraphBreak

还应分清集合上的两类条件。容量定义中的“\(u\geq1\) 几乎处处于某个开邻域”只涉及 Sobolev 等价类；“某个代表在 \(E\) 上逐点大于等于 \(1\)”涉及额外选择。下面先构造代表，再证明可对这些代表作点值容量估计，不倒过来用点值解释原定义。

\subsection{拟连续代表}
\label{b1:subsec:260202}

\begin{definition}\label{b1:def:sobolev-quasicontinuous}
设 \(v\colon\mathbb R^d\rightarrow\mathbb R\) 为有限值函数。若对每个 \(\varepsilon>0\)，都存在开集 \(G\)，使
\[
 C_p(G)<\varepsilon,\quad
 v|_{\mathbb R^d\setminus G}\ \text{连续},
\]
则称 \(v\) 为 \(p\)-\term[nilianxu]{拟连续}[quasicontinuous]函数。这里的连续性是补集上的相对连续性。若还有 \(v=u\) 几乎处处，其中 \(u\in W^{1,p}(\mathbb R^d)\)，则称 \(v\) 为 \(u\) 的拟连续代表。
\end{definition}

本节在容量零集上可以补以有限值，例如补零，不要求那些点具有由局部平均恢复的值。连续函数当然拟连续，只需取 \(G=\varnothing\)。另一个直接事实是：有限个拟连续函数的线性组合仍拟连续。为每个函数选取足够小的坏开集，取它们的并；在并集外，各项均连续，容量则由各项容量之和控制。

\begin{theorem}\label{b1:thm:sobolev-quasicontinuous-representative}
每个 \(u\in W^{1,p}(\mathbb R^d)\) 都有有限值 Borel 拟连续代表 \(\widetilde u\)。而且可以选取
\(\varphi_j\in C_c^\infty(\mathbb R^d)\)，使 \(\varphi_j\to u\) 于 \(W^{1,p}\)，并使这列函数在拟处处收敛到 \(\widetilde u\)。对每个 \(\varepsilon>0\)，还能删去一个容量小于 \(\varepsilon\) 的开集，使同一列函数在其补集上一致收敛。
\end{theorem}
\begin{proof}
由\cref{b1:thm:whole-space-sobolev-density}及范数等价性，选择
\[
 E_p(\varphi_j-u)^{1/p}\leq2^{-3j-2}.
\]
能量的 \(p\) 次方根满足三角不等式，因此
\[
 E_p(\varphi_{j+1}-\varphi_j)^{1/p}\leq2^{-3j}.
\]
定义坏开集
\[
 A_j=\{\,x:|\varphi_{j+1}(x)-\varphi_j(x)|>2^{-j}\,\},
 \quad G_N=\bigcup_{j\geq N}A_j.
\]
函数之差连续，故这些集合确实是开集。由\cref{b1:lem:continuous-capacity-level-set}和容量的次可加性，
\begin{equation}\label{b1:eq:quasicontinuous-bad-open-sets}
 C_p(G_N)\leq\sum_{j\geq N}C_p(A_j)
 \leq\sum_{j\geq N}2^{jp}E_p(\varphi_{j+1}-\varphi_j)
 \leq\sum_{j\geq N}2^{-2jp}.
\end{equation}
右端趋于零。令 \(Z=\bigcap_NG_N\)，由单调性得到 \(C_p(Z)=0\)。

固定 \(N\)。在闭集 \(F_N=\mathbb R^d\setminus G_N\) 上，对 \(k>j\geq N\)，
\[
 |\varphi_k(x)-\varphi_j(x)|
 \leq\sum_{i=j}^{k-1}|\varphi_{i+1}(x)-\varphi_i(x)|
 \leq\sum_{i=j}^{k-1}2^{-i}\leq2^{1-j}.
\]
所以这列连续函数在 \(F_N\) 上一致为 Cauchy 列，其极限有限且在 \(F_N\) 上连续。不同 \(F_N\) 上的极限在重叠处相同，因为它们来自同一列逐点极限。又有
\(\bigcup_NF_N=\mathbb R^d\setminus Z\)，故可定义
\[
 \widetilde u(x)=
 \begin{cases}
 \displaystyle\lim_{j\to\infty}\varphi_j(x),&x\notin Z,\\
 0,&x\in Z.
 \end{cases}
\]
集合 \(Z\) 为 Borel 集，在其补集上取 Borel 函数列的有限极限，再在 \(Z\) 上补零，所得函数仍为 Borel 函数。\eqref{b1:eq:quasicontinuous-bad-open-sets}及 \(F_N\) 上的一致收敛证明了拟连续性与定理最后一项。

还须核实这是原来 Sobolev 类的代表，而非仅仅某个点态极限。上述逼近速度给出
\[
 \sum_j\|\varphi_j-u\|_p^p<\infty.
\]
对非负项积分求和，得到
\(\sum_j|\varphi_j(x)-u(x)|^p<\infty\) 几乎处处，因而
\(\varphi_j(x)\to u(x)\) 几乎处处。另一方面，\(C_p(Z)=0\) 已由测度下界推出 \(m_d(Z)=0\)。在两个满测度集合的交上，点态极限只能是 \(u\)，所以 \(\widetilde u=u\) 几乎处处。
\end{proof}

这里的坏开集按容量控制，集合 \(F_N\) 不必有内点，也不必包含每个点的邻域。因此，补集上的一致极限只能直接给出相对连续性，不能将结论改写为通常连续性。光滑逼近的速度也不可省略：普通范数收敛只说明相邻差趋零，而证明需要加权后的能量级数可和。

\subsection{拟连续代表的唯一性}
\label{b1:subsec:260203}

先把连续函数的容量估计推广到拟连续代表。此时超水平集未必开，需要借助定义中的坏开集，把点态条件放入一个真正的开邻域。

\begin{proposition}\label{b1:prop:quasicontinuous-capacity-level-set}
设 \(v\) 是 \(u\in W^{1,p}(\mathbb R^d)\) 的拟连续代表。对每个 \(t>0\)，有
\begin{equation}\label{b1:eq:quasicontinuous-capacity-level-set}
 C_p\bigl(\{\,x:|v(x)|>t\,\}\bigr)
 \leq t^{-p}E_p(u).
\end{equation}
\end{proposition}
\begin{proof}
固定 \(\varepsilon>0\) 及 \(0<a<t\)。选取开集 \(G\)，使
\(C_p(G)<\varepsilon\)，且 \(v\) 在闭集 \(F=G^c\) 上连续。相对开集
\(\{\,x\in F:|v(x)|>a\,\}\) 可以写成 \(O\cap F\)，其中 \(O\) 在全空间开。于是
\[
 \{\,|v|>t\,\}\subseteq O\cup G,\quad
 |v|>a\ \text{于 }O\setminus G.
\]
由开集容量定义，取实值 \(w_G\in W^{1,p}\)，使
\[
 0\leq w_G\leq1,\quad
 w_G=1\ \text{几乎处处于 }G,\quad
 E_p(w_G)<C_p(G)+\varepsilon<2\varepsilon.
\]
截断到 \([0,1]\) 不增能量，故前两项可以同时要求。

在 Sobolev 等价类中定义
\[
 z=\max\{\,|u|/a,w_G\,\}.
\]
由\cref{b1:lem:sobolev-lattice-energy}，
\[
 z\in W^{1,p},\quad
 E_p(z)\leq a^{-p}E_p(u)+E_p(w_G)
 <a^{-p}E_p(u)+2\varepsilon.
\]
在 \(G\) 上，函数 \(w_G\) 几乎处处等于 \(1\)；在 \(O\setminus G\) 上，由 \(u=v\) 几乎处处可知 \(|u|/a>1\) 几乎处处。因此 \(z\) 是开集 \(O\cup G\) 的容许函数，进而
\[
 C_p\bigl(\{\,|v|>t\,\}\bigr)
 \leq C_p(O\cup G)
 \leq a^{-p}E_p(u)+2\varepsilon.
\]
先令 \(\varepsilon\downarrow0\)，再令 \(a\uparrow t\)，得到所需估计。
\end{proof}

估计中的能量由等价类 \(u\) 决定，左端的超水平集却由指定代表 \(v\) 决定。把两者连接起来的是拟连续性；如果给任意几乎处处代表套用同一估计，就会错误地认定所有 Lebesgue 零测集都具有零容量。

\begin{theorem}\label{b1:thm:quasicontinuous-uniqueness}
同一 Sobolev 类的两个有限值拟连续代表拟处处相等。反过来，在容量零集上任意改变一个拟连续代表的有限点值，所得函数仍是同一类的拟连续代表。
\end{theorem}
\begin{proof}
设 \(v,w\) 为同一类的拟连续代表。它们的差拟连续，并且几乎处处为零，故由前一命题，对每个正整数 \(j\)，
\[
 C_p\bigl(\{\,|v-w|>1/j\,\}\bigr)
 \leq j^p E_p(0)=0.
\]
取这些集合的可数并，得到 \(C_p(\{v\ne w\})=0\)。

再设 \(N\) 容量为零，函数 \(v'\) 在 \(N^c\) 上与 \(v\) 相同，在 \(N\) 上仍取有限值。给定 \(\varepsilon>0\)，先删去开集 \(G\)，使
\(C_p(G)<\varepsilon/2\)，且 \(v|_{G^c}\) 连续；再由外正则性选开集 \(H\supseteq N\)，使 \(C_p(H)<\varepsilon/2\)。在
\((G\cup H)^c\) 上，两个函数相同，故 \(v'\) 也连续，而
\(C_p(G\cup H)<\varepsilon\)。这就证明拟连续性。又因容量零集为 Lebesgue 零测集，且 Lebesgue 测度完备，任意这样的有限值改动仍可测，且不改变 Sobolev 类。
\end{proof}

因此可以用 \(\widetilde u\) 表示 Sobolev 类的拟连续代表，只要涉及的点态等式均在拟处处意义下理解。\symidx{\(\widetilde u\)：Sobolev 函数的拟连续代表，拟处处唯一}%
这个记号没有把某个光滑逼近序列或零容量集上的补值固定为额外结构。

\subsection{范数收敛与拟处处收敛}
\label{b1:subsec:260204}

拟连续代表的选择还与 Sobolev 范数中的极限相容。这比单纯的 \(L^p\) 几乎处处抽列结论更细，所用超水平估计则必须是刚刚证明的拟连续版本。

\begin{theorem}\label{b1:thm:sobolev-quasi-everywhere-subsequence}
若 \(u_n\to u\) 于 \(W^{1,p}(\mathbb R^d)\)，为各项及极限任取有限值拟连续代表 \(\widetilde u_n,\widetilde u\)，则存在子列
\(\widetilde u_{n_j}\) 拟处处收敛到 \(\widetilde u\)。对每个 \(\varepsilon>0\)，还能删去一个容量小于 \(\varepsilon\) 的开集，使该子列在补集上一致收敛。
\end{theorem}
\begin{proof}
抽取子列，使
\[
 E_p(u_{n_j}-u)\leq2^{-2jp}.
\]
差函数的拟连续代表可取
\(\widetilde u_{n_j}-\widetilde u\)。由超水平估计，
\[
 C_p\bigl(\{\,|\widetilde u_{n_j}-\widetilde u|>2^{-j}\,\}\bigr)
 \leq2^{-jp}.
\]
外正则性允许把每个超水平集包含在开集 \(O_j\) 中，并要求
\(C_p(O_j)<2^{1-jp}\)。令
\[
 V_N=\bigcup_{j\geq N}O_j,\quad Z=\bigcap_NV_N.
\]
则 \(C_p(V_N)\leq2\sum_{j\geq N}2^{-jp}\to0\)，因而 \(C_p(Z)=0\)。在 \(V_N^c\) 上，对每个 \(j\geq N\) 都有
\[
 |\widetilde u_{n_j}(x)-\widetilde u(x)|\leq2^{-j}.
\]
这同时证明补集上的一致收敛和 \(Z^c\) 上的逐点收敛。给定 \(\varepsilon\) 后取足够大的 \(N\)，即得最后一项。
\end{proof}

对整个收敛序列，虽未必有逐点收敛，上述命题仍给出每个固定 \(t>0\) 的定量结论：
\[
 C_p\bigl(\{\,|\widetilde u_n-\widetilde u|>t\,\}\bigr)
 \leq t^{-p}E_p(u_n-u)\longrightarrow0.
\]
这里不能用任意的可测代表替换拟连续代表，也不能只假设零阶 \(L^p\) 收敛而删去梯度控制。

\ProseParagraphBreak

再与\secref{b1:sec:1404}的精确代表 \(u^*\) 比较。那里的定义由缩小球的平均值决定，与可测代表的任意改动无关；本节的 \(\widetilde u\) 则由容量和连续逼近确定，只在拟处处意义下唯一。由于两者均与 \(u\) 几乎处处相等，已有
\[
 \widetilde u=u^*\quad\text{几乎处处}.
\]
若 \(\widetilde u\) 在某点 \(x\) 通常连续，则对充分小的球，
\[
 \frac1{m_d\bigl(B(x,r)\bigr)}
 \int_{B(x,r)}|u(y)-\widetilde u(x)|\,\mathrm dy
 \leq\sup_{y\in B(x,r)}
          |\widetilde u(y)-\widetilde u(x)|\longrightarrow0.
\]
所以这个点具有由平均绝对偏差确定的值，且 \(u^*(x)=\widetilde u(x)\)。对于一般拟连续点，删去小容量集后的相对连续性本身并未估计球内落在该坏集中的积分。本节因而不把上述几乎处处等式升级为“每个拟连续代表拟处处都是 Lebesgue 点”；这需要进一步的容量微分估计，不能仅从普通 Lebesgue 点定理推出。

\ProseParagraphBreak

单点容量把两种连续性的边界具体化。若 \(p>d\)，由\cref{b1:prop:sobolev-point-capacity}，每个非空集合的容量都有同一个正下界。拟连续定义中把 \(\varepsilon\) 取得小于这个下界时，坏开集只能为空，故拟连续函数实际上通常连续。若 \(1<p\leq d\)，单点容量为零，次可加性说明 \(C_p(\mathbb Q^d)=0\)。于是 \(\mathbf1_{\mathbb Q^d}\) 是零 Sobolev 类的一个拟连续代表，却处处不连续；它与连续代表零函数恰好在一个容量零集上不同。这也说明拟处处唯一不能替换为逐点唯一。

\ProseParagraphBreak

复值情形可按实部、虚部分别构造代表，再把两组坏开集合并。两个分量的唯一性给出复代表的拟处处唯一；范数收敛抽列时先后为两个分量选子列，仍得到共同坏开集外的一致收敛。容量的容许函数依旧是实值非负函数，不因讨论复值代表而改变。

% !TeX root = ../../Book1.tex
\section{容量与 Hausdorff 测度}
\label{b1:sec:2603}

容量由函数能量定义，Hausdorff 测度由小集合覆盖定义。两者相遇的地方是球截断：半径为 \(r\) 的球可以用一个梯度能量约为 \(r^{d-p}\) 的函数包住。这个尺度解释了 \(d-p\) 为何成为临界维数，但它本身还不能处理临界测度有限、却不为零的集合。

Kinnunen 与 Martio 的《The Sobolev capacity on metric spaces》\cite[pp. 380--381, 4.13. Theorem and 4.15. Theorem]{KinnunenMartio1996Capacity}组织了容量与 Hausdorff 覆盖之间的两个方向。本节使用本书的经典弱梯度，分别补出临界测度有限时的凸组合论证，以及反方向中的 \(5r\) 覆盖步骤。除非另行限制，以下集合不必为 Borel 集；\(\mathcal H^s\) 与 \(C_p\) 均按外测度值理解。

% 实读 KinnunenMartio1996Capacity 官方全文第380--381页4.13--4.15及完整证明，
% 并核看第381页。该文的g是Hajlasz梯度，本节不将其当作弱梯度。
% 球覆盖方向参考4.13的组织；反向以本书球Poincare重写4.15，
% 并用13章5r引理展开原文援引Federer的最后一步。
% H^{d-p}有限/σ有限推出零容量是本节独立展开的经典Sobolev凸组合论证。
% 前置：26.1格能量与球截断；9章Mazur；13章5r；15章Hausdorff外测度；
% 18章径向积分；21章Sobolev弱子列；24章球Poincare。不用Frostman或一般可容性。

\subsection{维数估计}
\label{b1:subsec:260301}

先设 \(1<p<d\)，记 \(t=d-p>0\)。由\cref{b1:prop:sobolev-capacity-balls}，当 \(0<r\leq1\) 时，
\[
 C_p\bigl(B(x,r)\bigr)\leq C_{d,p}(r^d+r^{d-p})
 \leq2C_{d,p}r^t.
\]
它立即给出一个覆盖上界。

\begin{proposition}\label{b1:prop:capacity-hausdorff-upper}
设 \(1<p<d\)。存在仅依赖 \(d,p\) 的常数 \(C\)，使每个
\(E\subseteq\mathbb R^d\) 都满足
\[
 C_p(E)\leq C\mathcal H^{d-p}(E).
\]
\end{proposition}
\begin{proof}
令 \(t=d-p\)，并假设右端有限。由 Hausdorff 覆盖定义，对每个
\(0<\delta<1\) 及 \(\eta>0\)，可以用开球 \(B(x_i,r_i)\) 覆盖 \(E\)，使
\begin{equation}\label{b1:eq:critical-hausdorff-ball-cover}
 0<r_i\leq\delta,\quad
 \sum_i r_i^t\leq C_t\mathcal H^t(E)+\eta.
\end{equation}
具体地，先取直径不超过 \(\delta/4\) 的可数覆盖。直径 \(a_i>0\) 的非空成员，可放入以其中一点为中心、半径 \(2a_i\) 的开球；直径为零的成员至多是单点，为这些单点另选正半径，使半径的 \(t\) 次幂之和小于任意给定误差。再利用归一化常数 \(c_t>0\)，便得到所列球覆盖。这里使用了 \(t>0\)。

单调性、可数次可加性及球上界给出
\[
 C_p(E)\leq\sum_iC_p\bigl(B(x_i,r_i)\bigr)
 \leq C_{d,p}\sum_i r_i^t
 \leq C\mathcal H^t(E)+C_{d,p}\eta.
\]
令 \(\eta\downarrow0\) 即得结论。右端无穷时不需另证。
\end{proof}

这个命题已说明 \(\mathcal H^{d-p}(E)=0\) 足以推出零容量。若临界测度只是有限，上界却只给出有限容量；要进一步得到零，还必须利用覆盖半径能够趋于零。

\begin{theorem}\label{b1:thm:finite-critical-hausdorff-zero-capacity}
设 \(1<p<d\)。若 \(E\subseteq\mathbb R^d\) 可以写成可数个具有有限
\(\mathcal H^{d-p}\) 外测度的集合之并，则 \(C_p(E)=0\)。特别地，
\(\mathcal H^{d-p}(E)<\infty\) 时容量为零。
\end{theorem}
\begin{proof}
先设 \(\mathcal H^t(E)<\infty\)，其中 \(t=d-p\)。取 \(0<\delta_n<1\)、\(\delta_n\downarrow0\)，并用\eqref{b1:eq:critical-hausdorff-ball-cover}选取开球覆盖
\[
 E\subseteq G_n=\bigcup_i B(x_{n,i},r_{n,i}),\quad
 r_{n,i}\leq\delta_n,\quad
 \sum_i r_{n,i}^t\leq M.
\]
常数 \(M<\infty\) 与 \(n\) 无关，可把每次覆盖的额外误差取为 \(1\)。

固定 \(0\leq\chi\leq1\)、\(\chi\in C_c^\infty\bigl(B(0,2)\bigr)\)，使 \(\chi=1\) 于 \(B(0,1)\)，并令
\[
 \chi_{n,i}(x)=\chi\left(\frac{x-x_{n,i}}{r_{n,i}}\right).
\]
直接伸缩计算给出
\[
 E_p(\chi_{n,i})
 \leq C_{d,p}(r_{n,i}^d+r_{n,i}^{d-p}).
\]
我们需要这些截断的可数最大值，而不是它们的和；求和可能在高度上重复付出代价。对固定 \(n\)，置
\[
 v_{n,k}=\max_{1\leq i\leq k}\chi_{n,i},\quad
 u_n=\sup_{i\geq1}\chi_{n,i},\quad
 S_n=\bigcup_i B(x_{n,i},2r_{n,i}).
\]
由\cref{b1:lem:sobolev-lattice-energy}，有
\[
 E_p(v_{n,k})\leq\sum_iE_p(\chi_{n,i})\leq C_{d,p}M.
\]
又有 \(0\leq v_{n,k}\uparrow u_n\leq\mathbf1_{S_n}\)，并且
\begin{equation}\label{b1:eq:critical-cover-small-volume}
 m_d(S_n)\leq2^d\omega_d\sum_i r_{n,i}^d
 \leq2^d\omega_d\delta_n^p\sum_i r_{n,i}^{d-p}
 \leq C_dM\delta_n^p.
\end{equation}
所以 \(v_{n,k}\to u_n\) 于 \(L^p\)。由\cref{b1:cor:sobolev-weak-subsequence}取一个 Sobolev 弱收敛子列，其零阶极限只能为 \(u_n\)。能量是一个等价范数的 \(p\) 次幂，故弱下半连续，得到
\[
 u_n\in W^{1,p}(\mathbb R^d),\quad
 E_p(u_n)\leq C_{d,p}M,\quad
 u_n=1\ \text{几乎处处于 }G_n.
\]
最后一项也可从逐点最大值直接读出；作为 Sobolev 类使用时，只要求这个几乎处处条件。

由\eqref{b1:eq:critical-cover-small-volume}，
\(\|u_n\|_p^p\leq C_dM\delta_n^p\to0\)。再从能量有界列
\((u_n)\) 中抽取 Sobolev 弱收敛子列，零阶强极限把其弱极限确定为零。仍把该子列记作 \((u_n)\)，并保留相应的 \(G_n\)。

由\cref{b1:thm:mazur}，每个尾部的有限凸组合可以在 Sobolev 范数中任意接近零。于是存在
\[
 z_j=\sum_{n=j}^{N_j}a_{j,n}u_n,\quad
 a_{j,n}\geq0,\quad \sum_{n=j}^{N_j}a_{j,n}=1,\quad
 E_p(z_j)\longrightarrow0.
\]
有限交 \(U_j=\bigcap_{n=j}^{N_j}G_n\) 仍是包含 \(E\) 的开集。在这个交集上，所有参与组合的 \(u_n\) 几乎处处等于 \(1\)，所以
\(z_j=1\) 几乎处处于 \(U_j\)。因此
\[
 C_p(E)\leq C_p(U_j)\leq E_p(z_j)\longrightarrow0.
\]
此处必须用有限凸组合；无限多个开邻域的交未必还是开邻域。

最后，若 \(E=\bigcup_\ell E_\ell\)，每个 \(E_\ell\) 的临界外测度有限，就由刚证的结论及容量的可数次可加性得到 \(C_p(E)=0\)。
\end{proof}

这个论证不要求 \(E\) 紧，也不要求从 \(E\) 上取得一个 Frostman 测度。它控制的是开集 \(S_n\) 的体积，不是 \(u_n\) 的拓扑支集；当覆盖的中心稠密时，\(S_n\) 的闭包可能很大。

\subsection{零容量集}
\label{b1:subsec:260302}

反方向的目标不同：不是从覆盖造容许函数，而是用一组能量很小的容许函数造出一个函数，使它在给定集合每一点附近的平均值都发散。这迫使梯度能量在那些点周围聚集。

\begin{theorem}\label{b1:thm:zero-capacity-hausdorff-dimension}
设 \(1<p\leq d\)，且 \(E\subseteq\mathbb R^d\) 满足 \(C_p(E)=0\)。则对每个 \(s>d-p\)，都有
\[
 \mathcal H^s(E)=0.
\]
特别地，零容量集的 Hausdorff 维数至多为 \(d-p\)。
\end{theorem}
\begin{proof}
由容量为零，选取开集 \(G_j\supseteq E\) 和非负函数
\(u_j\in W^{1,p}(\mathbb R^d)\)，使
\[
 u_j\geq1\ \text{几乎处处于 }G_j,\quad
 E_p(u_j)^{1/p}\leq2^{-j}.
\]
能量范数的可和性和 Sobolev 完备性说明
\[
 w=\sum_{j=1}^{\infty}u_j\in W^{1,p}(\mathbb R^d).
\]
这里可取非负逐点和作为几乎处处代表：部分和在 \(L^p\) 中收敛，又单调递增，故逐点和有限几乎处处，并与范数极限相同。

任取 \(x\in E\) 及正整数 \(N\)。有限交
\(\bigcap_{j=1}^NG_j\) 是 \(x\) 的开邻域，其中 \(w\geq N\) 几乎处处。因此
\begin{equation}\label{b1:eq:zero-capacity-divergent-means}
 \lim_{r\downarrow0}w_{B(x,r)}=+\infty
 \quad(x\in E).
\end{equation}
每个固定球上的平均仍有限，因为 \(w\in L^p\subseteq L^1_{\mathrm{loc}}\)。这个发散结论由开邻域上的几乎处处不等式推出，不需要预先给 \(w\) 选拟连续点值。

令 \(\mu=|\nabla w|^p\,m_d\)，它是有限正 Borel 测度。固定 \(s>d-p\)。假设某个 \(x\in E\) 满足
\[
 \limsup_{r\downarrow0}\frac{\mu\bigl(B(x,r)\bigr)}{r^s}<\infty.
\]
则存在 \(R>0\) 和 \(A<\infty\)，使
\(\mu\bigl(B(x,r)\bigr)\leq Ar^s\) 对 \(0<r\leq R\) 均成立。记
\(r_i=2^{-i}R\)、\(B_i=B(x,r_i)\)。由 Hölder 不等式及\cref{b1:thm:ball-poincare-lp}，
\[
 \begin{aligned}
 |w_{B_{i+1}}-w_{B_i}|
 &\leq\frac1{m_d(B_{i+1})}\int_{B_{i+1}}|w-w_{B_i}|\\
 &\leq C_d\,m_d(B_i)^{-1/p}\|w-w_{B_i}\|_{L^p(B_i)}\\
 &\leq C_d r_i^{1-d/p}\|\nabla w\|_{L^p(B_i)}
 \leq C_dA^{1/p}r_i^{(p-d+s)/p}.
 \end{aligned}
\]
指数 \((p-d+s)/p\) 严格为正，所以右端对 \(i\) 可和。相邻平均值之差可和，说明 \((w_{B_i})\) 收敛到有限值，与\eqref{b1:eq:zero-capacity-divergent-means}矛盾。因此
\begin{equation}\label{b1:eq:capacity-zero-gradient-density}
 E\subseteq
 \left\{\,x:\limsup_{r\downarrow0}
               \frac{\mu\bigl(B(x,r)\bigr)}{r^s}=+\infty\,\right\}.
\end{equation}

最后把这个密度条件变成覆盖估计。固定 \(M>0\)、\(\delta>0\) 和 \(R_0>0\)。对每个 \(x\in E\cap B(0,R_0)\)，由\eqref{b1:eq:capacity-zero-gradient-density}，存在任意小的半径
\(0<r<\delta/10\)，使
\(\mu\bigl(B(x,r)\bigr)>Mr^s\)。把全部这样的开球组成一个族。由\cref{b1:lem:five-r-covering}，可选至多可数个两两不交的球 \(B_i=B(x_i,r_i)\)，使
\[
 E\cap B(0,R_0)\subseteq\bigcup_i5B_i,\quad
 \operatorname{diam}(5B_i)=10r_i<\delta.
\]
于是
\[
 \begin{aligned}
 \mathcal H^s_\delta\bigl(E\cap B(0,R_0)\bigr)
 &\leq c_s10^s\sum_i r_i^s\\
 &\leq\frac{c_s10^s}{M}\sum_i\mu(B_i)
 \leq\frac{c_s10^s}{M}\mu(\mathbb R^d).
 \end{aligned}
\]
此处测度只作用于所选开球，不要求 \(E\) 可测。先令
\(\delta\downarrow0\)，再令 \(M\to\infty\)，得到
\(\mathcal H^s(E\cap B(0,R_0))=0\)。最后对整数 \(R_0\) 取可数并，便有 \(\mathcal H^s(E)=0\)。
\end{proof}

证明中严格不等式 \(s>d-p\) 用在几何级数的收敛上。若直接取 \(s=d-p\)，相邻平均值的估计只剩一个不随尺度衰减的常数，不能据此推出平均值收敛。

\subsection{奇异集}
\label{b1:subsec:260303}

上一节产生的例外集，现在可以附加几何大小的说明。这里所说的异常，是某条拟处处结论的例外，不是所有通常不连续点的统称。

\begin{corollary}\label{b1:cor:capacity-exception-dimension}
设 \(1<p\leq d\)。同一 Sobolev 类的两个拟连续代表不相等的集合，对每个 \(s>d-p\) 都具有零 \(\mathcal H^s\) 测度。若一列 Sobolev 函数在范数中收敛，则可以选取子列，使其拟连续代表不逐点收敛到极限代表的集合也满足同一结论。
\end{corollary}
\begin{proof}
由\cref{b1:thm:quasicontinuous-uniqueness}，第一类集合容量为零。由\cref{b1:thm:sobolev-quasi-everywhere-subsequence}，第二类集合可以包含在一个容量零集中。分别应用前一定理即可。
\end{proof}

例如，在 \(\mathbb R^3\) 中取 \(p=2\)，一条有限线段的
\(\mathcal H^1\) 测度有限，故其 \(2\)-容量为零。这个结论正好落在临界维数 \(d-p=1\) 上，不能仅凭“维数较低”而跳过有限临界测度的论证。

更一般地，在 \(1<p<d\) 时，若
\(\dim_{\mathrm H}E<d-p\)，则 \(\mathcal H^{d-p}(E)=0\)，所以容量为零；若
\(\dim_{\mathrm H}E>d-p\)，容量便不可能为零。两种严格不等式给出清楚的判别，而维数恰为 \(d-p\) 时仍须查看更细的测度或覆盖信息。特别是，已证的充分条件包括临界测度 \(\sigma\)-有限，并不只包括临界测度为零。

这些结果没有把普通 Lebesgue 点定理的例外集改成容量零集。\secref{b1:sec:2602}只在已经选定的拟连续代表与容量例外集之间建立了关系；由此得到的维数结论也只适用于这些例外集。例如当 \(1<p\leq d\) 时，零类的代表 \(\mathbf1_{\mathbb Q^d}\) 拟连续，却通常处处不连续，其通常不连续点集是整个空间。它不与本推论冲突，因为与零代表不相等的集合只是 \(\mathbb Q^d\)。

\subsection{临界维数}
\label{b1:subsec:260304}

当 \(p=d>1\) 时，幂次 \(r^{d-p}\) 变为常数，普通伸缩截断不能反映小球容量趋零的速度。这时可以把高度下降分散到许多同心环带，得到一个对数尺度。

\begin{proposition}\label{b1:prop:critical-logarithmic-capacity}
设 \(d>1\)。存在仅依赖于 \(d\) 的常数 \(\kappa_d\)，使
\[
 C_d\bigl(B(x,r)\bigr)
 \leq\kappa_d\bigl(\log(1/r)\bigr)^{1-d},
 \quad 0<r\leq\tfrac12.
\]
这里左端 \(C_d(\,\cdot\,)\) 是指数为 \(d\) 的容量。
\end{proposition}
\begin{proof}
平移不改变能量，故只考虑球心为零。令 \(L=\log(1/r)\)，取径向函数
\[
 v_r(x)=
 \begin{cases}
 1,&|x|\leq r,\\
 L^{-1}\log(1/|x|),&r<|x|<1,\\
 0,&|x|\geq1.
 \end{cases}
\]
对固定 \(r>0\)，该函数为全空间 Lipschitz 函数，紧支集且在
\(B(0,r)\) 上等于 \(1\)，因而是容许函数。在中间环带上，
\(|\nabla v_r(x)|=1/(L|x|)\) 几乎处处；其余部分梯度为零。
由\cref{b1:prop:radial-shell-integration}，
\[
 \int_{\mathbb R^d}|\nabla v_r|^d
 =d\omega_d L^{-d}\int_r^1\frac{\mathrm d\rho}{\rho}
 =d\omega_d L^{1-d}.
\]
非齐次容量还包括零阶项，不能只估计这一梯度积分。作代换
\(t=\log(1/\rho)\)，得到
\[
 \begin{aligned}
 \int_{\mathbb R^d}|v_r|^d
 &=\omega_dr^d+d\omega_dL^{-d}
       \int_r^1\bigl(\log(1/\rho)\bigr)^d\rho^{d-1}\,\mathrm d\rho\\
 &=\omega_de^{-dL}
   +d\omega_dL^{-d}\int_0^L t^de^{-dt}\,\mathrm dt
 \leq A_dL^{-d}.
 \end{aligned}
\]
最后一步使用
\(\sup_{L\geq\log2}L^de^{-dL}<\infty\) 及
\(\int_0^\infty t^de^{-dt}\,\mathrm dt<\infty\)；指数衰减使这两个量有限。因此，当 \(L\geq\log2\) 时，
\[
 E_d(v_r)\leq d\omega_dL^{1-d}+A_dL^{-d}
 \leq\kappa_dL^{1-d}.
\]
以这个能量作为球容量的上界，即得结论。
\end{proof}

为记录相应的覆盖判据，记
\[
 \ell_d(r)=\bigl(\log(1/r)\bigr)^{1-d},
 \quad 0<r\leq\tfrac12.
\]
对 \(0<\delta\leq1/2\)，直接用可数开球覆盖定义对数覆盖量
\[
 \mathcal H^{\ell_d}_\delta(E)
 =\inf\left\{\,
    \sum_i\ell_d(r_i):
    E\subseteq\bigcup_i B(x_i,r_i),\quad 0<r_i\leq\delta
   \,\right\}.
\]
再令
\[
 \mathcal H^{\ell_d}(E)
 =\lim_{\delta\downarrow0}\mathcal H^{\ell_d}_\delta(E).
\]
这个极限由尺度单调性存在。本式以半径为代价的自变量，没有另加
\(c_s\) 归一化；它是对数代价下的 Hausdorff 覆盖构造，不是
\(\mathcal H^0\) 的另一种记法。

\begin{corollary}\label{b1:cor:logarithmic-cover-zero-capacity}
若 \(d>1\)，则对任意 \(E\subseteq\mathbb R^d\)，有
\[
 C_d(E)\leq\kappa_d\mathcal H^{\ell_d}(E).
\]
特别地，对数覆盖量为零的集合具有零 \(d\)-容量。
\end{corollary}
\begin{proof}
对任意半径不超过 \(\delta\leq1/2\) 的开球覆盖，用容量的次可加性和前一命题，得到
\[
 C_d(E)\leq\sum_iC_d\bigl(B(x_i,r_i)\bigr)
 \leq\kappa_d\sum_i\ell_d(r_i).
\]
先对覆盖取下确界，再令 \(\delta\downarrow0\) 即可。
\end{proof}

由于 \(\ell_d(r)\to0\)，单点的对数覆盖量为零；可数集也可逐点分配任意小的可和代价，所以仍为零。另一方面，\(\mathcal H^0\) 是计数测度，单点已有 \(\mathcal H^0\) 值 \(1\)。这说明在临界指数处，以零维测度替代对数尺度会丢掉信息。

零 \(d\)-容量仍由前面定理推出 \(\mathcal H^s(E)=0\) 对全部 \(s>0\) 成立，但“维数为零”本身没有估计这里的对数覆盖量。也不能把 \(p<d\) 时有限临界测度的凸组合证明原样搬来断言有限对数覆盖量必然为零容量：上面的对数截断延伸到固定大小的球，未使容许函数在一个体积趋零的区域外为零。本节在此只使用已经证明的零覆盖量判据。

% !TeX root = ../../Book1.tex
\section{细拓扑}
\label{b1:sec:2604}

拟连续性允许删去容量很小的开集，却没有把余集变成通常的开集。若希望直接谈论这些代表的连续性，就需要一种能够辨认容量大小的邻域。本节先从相对容量构造这种拓扑，再说明它与边界值问题的联系。相对容量的基本性质和拓扑公理将在这里证明；Cartan 性质、拟连续与细连续的联系以及非线性 Wiener 判别属于进一步的势论，均明确给出证明概要，不把它们当成前面 Sobolev 理论的直接推论。

“细拓扑”与“薄性”采用 Doob 的《经典位势论与概率位势论（上册）》中文目录中的名称。\footnote{科学出版社的\href{https://www.ecsponline.com/goods.php?id=6199}{官方书目}第XI章及第XI.2、第XIII.17节采用这些名称；这里仅据目录核对称谓，不以古典位势论的目录代替非线性 \(p\)-容量的定义或证明。}本节的“细”指一套特定的势论拓扑，并不是泛指任意比原拓扑更细的拓扑。

% 中文“细拓扑”“薄性”据科学出版社《经典位势论与概率位势论（上册）》
% 官方目录第XI章、XI.2及XIII.17；目录只作名称证据，不作非线性证明来源。
% 实读 Heinonen--Kilpeläinen--Martio 1989, AIF 39(2), 293--318：
% 293--295页引言，298--299页第3节定义及Theorem 3.2；
% 300--303页只读相关证明路线，不声称完整重建全部势论。
% 实读 Björn--Björn--Latvala 2018, J. Anal. Math. 135(1), 59--83：
% 62页Theorem 1.4，64页拟开/拟连续定义，71页Definitions 6.1--6.2，
% 75页第8节Theorem 1.4的证明；未读完其所调用的所有先前论文。
% 实读 Kilpeläinen--Malý 1994, Acta Math. 172(1), 137--161：
% 137--139页Theorems 1.1、1.3、1.6的陈述；
% 155--158页第5节Cartan性质、Perron比较与Wiener判别证明路线。
% 未完整读第4节非线性势估计的证明；本节将该估计明确列为未展开输入。
% 相对开集a.e.定义、格能量证明和离散尺度整理按本书前章工具重写。

\subsection{薄性}
\label{b1:subsec:260401}

本节仍取 \(1<p<\infty\)。全空间容量 \(C_p\) 同时计入函数和梯度的能量；在一个不断缩小的球内考察边界障碍时，更合适的是固定外边界、只计梯度的相对容量。两个对象控制相同的局部零容量现象，但尺度公式不同，不能混用。

\begin{definition}\label{b1:def:relative-variational-capacity}
设 \(D\subset\mathbb R^d\) 为有界开集。给定开集 \(G\subset D\)，定义它相对于 \(D\) 的 \term[xiangdui-bianfen-rongliang]{相对变分容量}[relative variational capacity]为
\[
 \operatorname{cap}_p(G,D)
 =\inf\left\{\int_D|\nabla v|^p\,\mathrm dx:
 v\in W_0^{1,p}(D),\
 v\ge1\ \text{几乎处处于 }G\right\}.
\]
空容许类的下确界为 \(+\infty\)。对任意 \(E\subset D\)，再令
\[
 \operatorname{cap}_p(E,D)
 =\inf\{\,\operatorname{cap}_p(G,D):
 E\subset G\subset D,\ G\text{开}\,\}.
\]
这里的 \(W_0^{1,p}(D)\) 仍是内部紧支集光滑函数的闭包，不预先用边界点值解释。
\end{definition}
\symidx{\(\operatorname{cap}_p(E,D)\)：相对于开集的纯梯度容量}

\begin{lemma}\label{b1:lem:relative-capacity-basic}
相对容量关于 \(E\) 单调且可数次可加，关于外域 \(D\) 反向单调。容许函数可以限制为 \(0\le v\le1\)。对紧集 \(K\Subset D\)，还有
\begin{equation}\label{b1:eq:relative-compact-capacity}
 \operatorname{cap}_p(K,D)
 =\inf\left\{\int_D|\nabla\varphi|^p\,\mathrm dx:
 \varphi\in C_c^\infty(D),\
 \varphi\ge1\text{ 于 }K\text{ 的某个开邻域}\right\}.
\end{equation}
\end{lemma}
\begin{proof}
先说明格运算不会破坏零边界闭包。以 \(C_c^\infty(D)\) 逼近 \(v\in W_0^{1,p}(D)\)，对逼近函数取正部或截断。所得函数有内部紧支集，仍可在 \(W^{1,p}\) 中用内部光滑函数逼近。截断的梯度界给有界性，零阶项则收敛到相应的截断；弱极限识别及 \(W_0^{1,p}(D)\) 的弱闭性说明极限仍在这个空间。由正部公式又得到最大值、最小值的闭包性质。于是前面的格能量论证同样给出
\[
 \int_D|\nabla(v\vee w)|^p
 +\int_D|\nabla(v\wedge w)|^p
 =\int_D|\nabla v|^p+\int_D|\nabla w|^p.
\]
取值截断到 \([0,1]\) 不增加能量，也不改变容许性。

集合单调性来自容许类的包含；外域增大时，利用\cref{b1:prop:w0-zero-extension}把容许函数补零即可。证明可数次可加性时，只须考虑容量之和有限的情形。为各集选取开覆盖及近乎最优的容许函数 \(v_j\in[0,1]\)，使梯度能量之和有限，并令 \(w_N=\max_{j\le N}v_j\)。格公式控制其梯度，\cref{b1:thm:zero-boundary-poincare}控制其零阶项；而
\[
 0\le w_N\uparrow w,\quad
 w^p\le\sum_j v_j^p\in L^1(D).
\]
故 \(w_N\to w\) 于 \(L^p(D)\)。抽取 Sobolev 弱子列，极限仍在 \(W_0^{1,p}(D)\)，且梯度能量不超过各项能量之和。它在开覆盖之并上几乎处处等于 \(1\)。最后让选取容许函数时的总误差趋零。

式\eqref{b1:eq:relative-compact-capacity}右边至少等于左边。反过来，选取开集 \(G\supset K\) 和在 \(G\) 上几乎处处等于 \(1\) 的容许函数 \(v\in[0,1]\)。取 \(\psi\in C_c^\infty(G)\)，使 \(0\le\psi\le1\) 且在 \(K\) 的邻域等于 \(1\)，再取 \(v_j\in C_c^\infty(D)\) 趋于 \(v\)。函数
\[
 \varphi_j=\psi+(1-\psi)v_j
\]
在该邻域等于 \(1\)，并且
\[
 \varphi_j-v=(1-\psi)(v_j-v),
\]
因为 \(\psi(1-v)=0\) 几乎处处。光滑乘法的有界性给出 \(\varphi_j\to v\) 于 \(W^{1,p}(D)\)，所以光滑容许函数的能量可逼近原下确界。
\end{proof}

这也说明本定义与经典势论先在紧集上定义容量的口径一致。具体地，取开集 \(G\) 的递增紧穷竭 \(K_j\)，使 \(\bigcup_j\operatorname{int}K_j=G\)。若这些紧集的容量有统一上界，对其光滑容许函数取弱极限，每个固定 \(K_j\) 上的几乎处处约束都保留下来，因而极限在 \(G\) 上几乎处处至少为 \(1\)。这就证明
\(\operatorname{cap}_p(G,D)=\sup_{K\Subset G}\operatorname{cap}_p(K,D)\)；最后再对任意集取开外正则化。此处不需要先定义某个 Sobolev 代表在任意集合上的逐点约束。

\begin{proposition}\label{b1:prop:relative-capacity-scale}
平移、伸缩给出
\[
 \operatorname{cap}_p(x+rE,x+rD)
 =r^{d-p}\operatorname{cap}_p(E,D),\quad r>0.
\]
存在仅依赖 \(d,p\) 的正常数 \(c,C\)，使
\begin{equation}\label{b1:eq:relative-ball-capacity}
 cr^{d-p}\le
 \operatorname{cap}_p\bigl(B(x,r),B(x,2r)\bigr)
 \le Cr^{d-p}.
\end{equation}
此外，对任意 \(F\subset B(x,r)\)，有
\begin{equation}\label{b1:eq:relative-outer-ball-comparison}
 \operatorname{cap}_p\bigl(F,B(x,4r)\bigr)
 \le\operatorname{cap}_p\bigl(F,B(x,2r)\bigr)
 \le C\operatorname{cap}_p\bigl(F,B(x,4r)\bigr).
\end{equation}
\end{proposition}
\begin{proof}
作变量代换 \(y=x+rz\)，梯度产生因子 \(r^{-1}\)，体积产生因子 \(r^d\)；该变换及其逆保持零边界闭包和开邻域约束，所以先对开集、再对任意集得到精确的尺度公式。

固定一个在单位球上等于 \(1\)、支集含于二倍球的光滑截断，伸缩后得到球对的上界。若 \(v\) 是内球的容许函数，则零边界 Poincaré 不等式给出
\[
 \omega_d r^d
 \le\int_{B(x,2r)}|v|^p
 \le(4r)^p\int_{B(x,2r)}|\nabla v|^p,
\]
从而得到下界。

最后一个式子的左侧来自外域单调性。证明右侧时，取在 \(B(x,r)\) 上等于 \(1\)、支集含于 \(B(x,2r)\) 的截断 \(\eta\)，其中 \(|\nabla\eta|\le C/r\)。将 \(B(x,4r)\) 中的容许函数 \(v\) 乘以 \(\eta\)，便得到较小外球中的容许函数。乘积法则和零边界 Poincaré 不等式给出
\[
 \int|\nabla(\eta v)|^p
 \le C\int|\nabla v|^p+Cr^{-p}\int|v|^p
 \le C'\int|\nabla v|^p.
\]
若原约束在开集 \(G\supset F\) 上成立，乘积的约束就在 \(G\cap B(x,r)\) 上成立，仍是 \(F\) 的开邻域。取下确界即得结论。
\end{proof}

特别地，\(p=d\) 时球对容量有与半径无关的正上下界。这不与全空间容量 \(C_d(B(x,r))\to0\) 矛盾：相对容量要求函数在距离中心仅为 \(2r\) 的地方已经满足零边界条件。

\begin{definition}\label{b1:def:p-thin-set}
给定 \(E\subset\mathbb R^d\) 和 \(x\in\mathbb R^d\)，令 \(r_j=2^{-j}\)。若
\begin{equation}\label{b1:eq:p-wiener-series}
 \sum_{j=0}^{\infty}
 \left(
 \frac{\operatorname{cap}_p\bigl(E\cap B(x,r_j),B(x,2r_j)\bigr)}
 {r_j^{d-p}}
 \right)^{1/(p-1)}<\infty,
\end{equation}
则称 \(E\) 在 \(x\) 处具有 \(p\)-\term[boxing]{薄性}[thinness]，或称 \(E\) 在 \(x\) 处是 \(p\)-薄的；级数发散时则称它在该点不是 \(p\)-薄的。
\end{definition}

分母也可以换成球对容量，因\eqref{b1:eq:relative-ball-capacity}只会改变正常数。有限个大尺度项不影响收敛性。外球的倍数也不是定义的本质：\eqref{b1:eq:relative-outer-ball-comparison}及其相同的截断证明允许把倍数 \(2\) 换成任意固定的 \(\lambda>1\)。

文献通常在连续的对数尺度上写 Wiener 积分；本节采用二进级数，使任意集合 \(E\) 都可直接进入定义。两种尺度判据的比较只需前述估计：当 \(r/2\le t\le r\) 时，利用内集单调性、外球比较以及 \(t^{d-p}\) 与 \(r^{d-p}\) 的可比性，中间尺度的归一化容量被相邻两个二进尺度的量从上下控制。因此逐个对数尺度求和与通常的 Wiener 积分给出同一收敛判据，也不依赖起始半径的选择。相对容量及这一薄性条件可与 Heinonen、Kilpeläinen 和 Martio 的《Fine topology and quasilinear elliptic equations》\cite[Section 3]{Heinonen1989Fine}对读。

\subsection{细连续性}
\label{b1:subsec:260402}

\begin{definition}\label{b1:def:p-fine-topology}
设 \(U\subset\mathbb R^d\)。若 \(\mathbb R^d\setminus U\) 在每个 \(x\in U\) 处都是 \(p\)-薄的，则称 \(U\) 为 \(p\)-细开集。所有这样的集合组成的拓扑称为 \(p\)-\term[xituopu]{细拓扑}[fine topology]。
\end{definition}

这里先给出了开集的判据，仍须验证它确实定义一个拓扑。还要留意量词：补集只在一个点 \(x\) 处是薄的，并不根据这一定义就成为 \(x\) 的细邻域；定义要求对细开集中的每个点检查薄性。二者之间的联系将在最后一小节作为势论定理陈述。

\begin{proposition}\label{b1:prop:p-fine-topology-basic}
细开集组成一个比欧氏拓扑更细的拓扑。若 \(C_p(N)=0\)，则 \(\mathbb R^d\setminus N\) 为细开集。若 \(p>d\)，细拓扑恰好等于欧氏拓扑。
\end{proposition}
\begin{proof}
薄性对取子集保持。相对容量的次可加性与
\[
 (a+b)^q\le\max(1,2^{q-1})(a^q+b^q),\quad
 a,b\ge0,\quad q=\frac1{p-1},
\]
又说明有限个在 \(x\) 处薄的集合之并仍在 \(x\) 处薄。因此有限个细开集的交是细开的。若 \(x\in\bigcup_\lambda U_\lambda\)，选取一个包含 \(x\) 的 \(U_\lambda\)；并集的补集包含于这个 \(U_\lambda\) 的补集，所以也在 \(x\) 处薄。这证明任意并的封闭性，空集与全空间则直接满足条件。普通开集的补集在每个内点附近为空，故一定细开。

若 \(C_p(N)=0\)，固定 \(B(x,r)\)。将全空间中覆盖 \(N\) 的容许函数乘以在该球上等于 \(1\)、支集含于 \(B(x,2r)\) 的截断，得到
\[
 \operatorname{cap}_p\bigl(N\cap B(x,r),B(x,2r)\bigr)
 \le C_{r,p}\,C_p(N)=0.
\]
于是 \(N\) 在每一点都薄，结论成立。

最后设 \(p>d\)，并令 \(y\in B(x,r)\)。对 \(B(x,2r)\) 中覆盖单点 \(y\) 的任一容许函数，先补零到全空间。\cref{b1:thm:morrey-whole-space}给出其连续代表；它在容许开邻域上几乎处处至少为 \(1\)，故在 \(y\) 处也至少为 \(1\)。取 \(z=x+3re_1\)，补零代表在 \(z\) 处为零，因而
\[
 1\le |v(y)-v(z)|
 \le C(4r)^{1-d/p}\|\nabla v\|_{L^p}.
\]
所以相对单点容量至少为 \(cr^{d-p}\)，常数对 \(y\in B(x,r)\) 一致。若 \(x\in\overline E\)，每个小球都含有 \(E\) 的一点，\eqref{b1:eq:p-wiener-series}的每项便有统一正下界，不可能收敛。反之，若 \(x\notin\overline E\)，小尺度项最终为零。于是这个范围内，“在 \(x\) 处薄”等价于“在 \(x\) 的某个普通邻域中为空”，细开集恰好就是普通开集。
\end{proof}

当 \(1<p\le d\) 时则不同。由\cref{b1:sec:2601}的单点容量计算，\(\mathbb Q^d\) 的容量为零，因此 \(\mathbb R^d\setminus\mathbb Q^d\) 是细开集；它却不是欧氏开集，因为每个普通球都含有有理点。细拓扑所忽略的是某种容量很小的障碍，而不是距离上很远的点。

\begin{definition}\label{b1:def:p-fine-continuity}
设 \(v\colon\mathbb R^d\to\mathbb R\)。若对每个 \(\varepsilon>0\)，都存在含 \(x\) 的细开集 \(V\)，使
\[
 |v(y)-v(x)|<\varepsilon\quad(y\in V),
\]
则称 \(v\) 在 \(x\) 处 \(p\)-\term[xilianxu]{细连续}[finely continuous]。这就是定义域取细拓扑、值域取通常拓扑时的点态连续性。
\end{definition}

普通连续性蕴含细连续性。反过来没有这样的普遍蕴含；而拟连续性与细连续性的关系，也不是仅由“细拓扑更细”这一事实决定的。

\begin{theorem}\label{b1:thm:quasicontinuous-fine-continuity}
设 \(1<p<\infty\)。每个有限实值的 \(p\)-拟连续函数都在 \(p\)-拟处处细连续。特别地，\cref{b1:thm:sobolev-quasicontinuous-representative}所构造的 Sobolev 代表具有这一性质。
\end{theorem}
\begin{proofsketch}
需要的深层输入是以下集合刻画：若一个集合在加入容量任意小的适当开集后总能成为开集，则它可写成一个细开集与一个容量零集的并。这是非线性势论中的结果，其证明依赖 Cartan 型分离、Choquet 性质和相应的容量例外集控制；这些步骤不在本节展开。Björn、Björn 和 Latvala 的《The Cartan, Choquet and Kellogg properties for the fine topology on metric spaces》\cite[Theorem 1.4, Section 8]{Bjorn2018Cartan}给出该刻画与本定理的联系，这里仅使用其欧氏空间情形。

有了该输入，最后一步是直接的。拟连续函数 \(v\) 的每个超水平集 \(\{v>a\}\) 和次水平集 \(\{v<a\}\)，在加入使 \(v\) 于其补集连续的坏开集后，都成为开集。对所有有理数 \(a\) 应用集合刻画，合并所得的可数个容量零集为 \(N\)。在 \(N\) 外，这些水平集都由相应的细开集给出。又因 \(\mathbb R^d\setminus N\) 细开，围绕 \(v(x)\) 选取两个有理阈值，就得到细连续性所需的邻域。

因此，未展开的关键在于上述集合刻画，而不在有理阈值这一步。前一节在坏开集外的一致逼近本身，没有证明那些补集在其各点附近都是细邻域。
\end{proofsketch}

\subsection{Wiener 型思想}
\label{b1:subsec:260403}

薄性之所以采用指数 \(1/(p-1)\)，最终与 \(p\) 次梯度能量的非线性边界行为有关。为说明这点，先写出相应的方程。

\begin{definition}\label{b1:def:p-harmonic}
设 \(\Omega\subset\mathbb R^d\) 开。给定连续函数 \(u\in W_{\mathrm{loc}}^{1,p}(\Omega)\)，若
\[
 \int_\Omega |\nabla u|^{p-2}\nabla u\cdot\nabla\varphi\,\mathrm dx=0
 \quad(\varphi\in C_c^\infty(\Omega)),
\]
则称 \(u\) 为 \term[p-tiaohe-hanshu]{\(p\)-调和函数}[p-harmonic function]。当梯度为零时，向量 \( |\nabla u|^{p-2}\nabla u\) 约定为零。
\end{definition}

形式上，该方程写成 \(\Delta_pu=0\)，其中
\[
 \Delta_pu=\operatorname{div}\bigl(|\nabla u|^{p-2}\nabla u\bigr).
\]
\symidx{\(\Delta_pu\)：由 \(p\) 次梯度能量产生的非线性算子}
当 \(p=2\) 时，它是通常的 Laplace 方程；一般 \(p\) 时不能利用解的叠加。

设 \(\Omega\) 有界，且 \(g\in W^{1,p}(\Omega)\cap C(\overline\Omega)\)。在仿射类 \(g+W_0^{1,p}(\Omega)\) 上最小化 \(\int_\Omega|\nabla u|^p\)，确实有唯一极小元：对 \(u-g\) 用零边界 Poincaré 不等式，便由梯度界得到 Sobolev 有界性；弱子列、仿射类弱闭性和凸能量的弱下半连续性给出存在。严格凸性先给极小元的梯度相等，再由零边界 Poincaré 不等式得到函数相等。沿内部测试函数作变分，则得到上面的弱方程。这部分与\cref{b1:thm:sobolev-variational-minimum}使用同一直接法。

但极小元具有内部连续代表，尤其在 \(p\le d\) 时，不是 Morrey 不等式的推论，而是 \(p\)-调和方程的内正则性定理。以下边界判别把这项正则性作为经典非线性势论的输入，并用 \(u_g\) 表示该连续代表。

\begin{definition}\label{b1:def:p-dirichlet-regular-point}
设 \(\Omega\) 为有界开集，\(x_0\in\partial\Omega\)。若对每个 \(g\in W^{1,p}(\Omega)\cap C(\overline\Omega)\)，上述变分解都满足
\[
 \lim_{\substack{x\to x_0\\x\in\Omega}}u_g(x)=g(x_0),
\]
则称 \(x_0\) 为这个 \(p\)-Dirichlet 问题的\term[zhengze-bianjiedian]{正则边界点}[regular boundary point]。
\end{definition}

边界值在这里不是任意指定的可测代表值。闭包条件给出变分意义的边界约束，而正则边界点要求这一约束能够由内部连续解的极限恢复。

\term[wiener-panbie-zhunze]{Wiener 判别准则}[Wiener criterion]把这个分析问题转化为补集在每个小尺度上的容量问题。
\begin{theorem}[Wiener]\label{b1:thm:p-wiener-criterion}
设 \(d\ge2\)、\(1<p\le d\)，\(\Omega\subset\mathbb R^d\) 为有界开集，且 \(x_0\in\partial\Omega\)。则 \(x_0\) 是正则边界点，当且仅当 \(\mathbb R^d\setminus\Omega\) 在 \(x_0\) 处不是 \(p\)-薄的，即\eqref{b1:eq:p-wiener-series}中取 \(E=\mathbb R^d\setminus\Omega\) 所得级数发散。
\end{theorem}
\begin{proofsketch}
充分性需要把补集的相对容量转化为变分解在相邻尺度间的边界振荡衰减，再迭代这些估计。发散条件恰好使未被边界约束消除的振荡趋于零。这里需要方程的比较原理和容量控制的边界衰减估计，不能用普通球均值 Poincaré 不等式直接代替。

必要性更深：当补集薄时，借助相对容量的极小函数、非线性势的点态估计及 Cartan 性质，构造在该点不能恢复指定边值的解。Kilpeläinen 与 Malý 的《The Wiener test and potential estimates for quasilinear elliptic equations》\cite[THEOREM 1.1, Section 5]{Kilpelainen1994Wiener}给出上述完整范围，并说明 Cartan 性质、比较方法和 Wiener 判别之间的联系。本节不展开支撑必要性的非线性势估计，也不另行建立 Perron 方法及其与变分解的对应。定理在这里用于指出容量理论通向边界正则性的方向，而不是宣称这些势论工具已在前文证明。
\end{proofsketch}

这个判据可以立即解释两种相反的几何情形。若存在固定 \(a>0\)，使每个充分小的 \(B(x_0,r)\setminus\Omega\) 都包含一个半径为 \(ar\) 的球，那么任何容许函数在该小球上至少为 \(1\)。对外球中的容许函数用零边界 Poincaré 不等式，得到
\[
 \operatorname{cap}_p\bigl(B(x_0,r)\setminus\Omega,B(x_0,2r)\bigr)
 \ge c_a r^{d-p}.
\]
故 Wiener 级数发散。Lipschitz 图表的外侧含有这种按比例缩小的球，因此它的边界点符合该判据。相反，在穿孔球 \(\Omega=B(0,1)\setminus\{0\}\) 中，\(1<p\le d\) 时内边界点 \(0\) 附近的补集只有容量为零的单点，级数收敛，于是该点不正则。这说明一个拓扑边界点未必足以对变分解施加可见的点态边值。

\subsection{势论联系}
\label{b1:subsec:260404}

线性势论以调和函数及超调和函数组织比较原理；非线性情形保留比较的思想，但不能把超调和函数都当作局部 Sobolev 函数。

\begin{definition}\label{b1:def:p-superharmonic}
设 \(\Omega\subset\mathbb R^d\) 开。给定下半连续、在一个稠密子集上有限的函数 \(v\colon\Omega\to(-\infty,+\infty]\)。若对每个满足 \(\overline D\Subset\Omega\) 的开集 \(D\)，以及每个在 \(D\) 上 \(p\)-调和且属于 \(C(\overline D)\) 的 \(h\)，都有
\[
 h\le v\text{ 于 }\partial D
 \quad\Longrightarrow\quad h\le v\text{ 于 }D,
\]
则称 \(v\) 为 \term[p-chaotiaohe-hanshu]{\(p\)-超调和函数}[p-superharmonic function]。
\end{definition}

定义允许取 \(+\infty\)，也不先要求 \(v\in W_{\mathrm{loc}}^{1,p}\)。因此不能未经论证就把每个这样的 \(v\) 代入弱方程并写出其梯度。超调和函数的截断正则性、下半连续代表与弱不等式的对应，都属于这套理论需要另外建立的内容。

\term[cartan-xingzhi]{Cartan 性质}[Cartan property]说明，薄性可以由一个势函数在该点与沿障碍集合的值之间的严格间隙检测。
\begin{theorem}[Cartan]\label{b1:thm:p-cartan-characterization}
设 \(d\ge2\)、\(1<p\le d\)，且 \(x_0\in\overline E\setminus E\)。集合 \(E\) 在 \(x_0\) 处是 \(p\)-薄的，当且仅当在 \(x_0\) 的某个普通邻域中存在非负 \(p\)-超调和函数 \(v\)，使
\[
 v(x_0)<
 \liminf_{\substack{x\to x_0\\x\in E}}v(x).
\]
相应地，对任意 \(U\ni x_0\)，\(U\) 包含 \(x_0\) 的一个细开邻域，当且仅当 \(\mathbb R^d\setminus U\) 在 \(x_0\) 处是 \(p\)-薄的。细拓扑也正是使所有局部 \(p\)-超调和函数连续的最弱拓扑，允许函数值为 \(+\infty\) 时在值域采用通常的扩充实数拓扑。
\end{theorem}
\begin{proofsketch}
由势函数的比较性质和容量估计，可以从严格间隙推出薄性。反向需要将小尺度的容量障碍组合成超调和势，同时控制它在 \(x_0\) 的值。在线性情形可以使用势的叠加；一般 \(p\) 时必须用非线性势估计与比较方法替代，不能把线性证明中的级数逐项照搬。

Kilpeläinen 与 Malý 的《The Wiener test and potential estimates for quasilinear elliptic equations》\cite[THEOREM 1.3, Section 5]{Kilpelainen1994Wiener}建立这种点态分离；由超调和函数定义的拓扑与薄性邻域的等价刻画，可参阅 Heinonen、Kilpeläinen 和 Martio 的《Fine topology and quasilinear elliptic equations》\cite[3.2. THEOREM]{Heinonen1989Fine}。这些结论包含了前文尚未证明的邻域提取步骤。若一个集合已经含有细开邻域，其补集薄只需单调性；从仅在一个点成立的薄性反过来提取细开邻域，才是此处需要势论的方向。
\end{proofsketch}

可以用一个公式看清这些深层估计增加了什么。给定正 Radon 测度 \(\sigma\)，相关的局部非线性势积分具有形式
\[
 \int_0^R
 \left(\frac{\sigma(B(x,r))}{r^{d-p}}\right)^{1/(p-1)}
 \frac{\mathrm dr}{r}.
\]
它把测度在所有小尺度上的集中程度累积起来。完整理论先为容量障碍构造极小函数和相应测度，再用这类势积分估计超调和函数的点值；由此才能控制 Cartan 分离和边界极限。本节没有证明该测度构造及点态估计，公式本身也不能代替它们。

到这里，三种“忽略异常”的方式便有了明确分工。几乎处处只识别积分等价类；拟处处用容量限制可忽略的点集；细拓扑则让容量所控制的局部障碍进入邻域结构。前两者的基本 Sobolev 结论已在本章建立，后者与非线性方程之间的深层联系在本节给出了准确入口。继续阅读时，可先沿 Cartan 性质与 Wiener 判别的两篇原始论文补足势估计和比较方法，再回看拟连续代表为何能在拟处处成为细连续函数。


\begin{chaptersummary}
容量所度量的是在集合邻域保持函数大小所需的能量，而不是单独在集合上任意填写点值的代价。开集上的几乎处处约束与开外正则化，使这一概念能先于代表的选择建立。次可加性随后把快速逼近中的坏集汇合起来，得到拟连续代表及其拟处处唯一性。

\ProseParagraphBreak

零容量集一定零测，反向一般不成立。小球缩放、临界对数过渡与球平均的累积，把这种区别同 Hausdorff 维数联系起来；等于临界维数时，还要考察比单一维数更细的覆盖信息。拟连续性给出了小容量开集外的相对连续性，并不自动说明所有点都通常连续，也不能单靠普通微分定理确定其容量异常集。

\ProseParagraphBreak

细拓扑让邻域概念也适应容量尺度。它与边界可达值、超解及非线性势估计之间有深层联系，本章在相应证明概要中列明了尚未展开的输入。下一章回到可以直接由分布导数测度描述的对象：当弱导数不再有可积密度，而只是一份有限测度时，便进入多维有界变差函数理论。
\end{chaptersummary}



容量的单点行为随指数跨过维数而改变，见\cref{b1:tab:visual-t14}。
% Source fingerprint: 88b089aaadd01ba58e73e9a153026309df6f8ac9e9e35b19f12fcd69aca8859f
% T14: raw TeX cells preserved; no numeric highlighting.
\begin{table}[htbp]
\centering\small\renewcommand{\arraystretch}{1.18}\setlength{\tabcolsep}{4pt}
\caption[容量与维数的三个指数区间]{容量与维数的三个指数区间。前三行讨论本章全空间非齐次容量。最后一行是第26.4节建立的相对容量；两种能量的尺度不同。}
\label{b1:tab:visual-t14}
\begin{tabularx}{\linewidth}{@{}>{\raggedright\arraybackslash}p{3.1cm}>{\raggedright\arraybackslash}p{4.7cm}>{\raggedright\arraybackslash}X@{}}
\toprule
指数 & 小球与单点 & 零测集与零容量集 \\
\midrule
\(1<p<d\) & \(C_p(B_r)\le C(r^d+r^{d-p})\)；单点容量零 & 单点给出非空零容量集；平面片可为零测而正容量 \\
\(p=d>1\) & 需用对数型截断证明单点容量零；上一行幂估计本身不够 & 可数集容量零，但不能把所有 Lebesgue 零测集都忽略 \\
\(p>d\) & 单点容量有统一正下界 & 每个非空集合容量为正；只有空集可为零容量集 \\
相对容量的球对 & \(\operatorname{cap}_p(B_r,B_{2r})\) 与 \(r^{d-p}\) 可比 & 纯梯度能量，不能与全空间非齐次容量混为同一数值 \\
\bottomrule
\end{tabularx}
\end{table}


\BookExercises


% \chapref{b1:ch:26}习题临时稿；不另设 BookExercises 入口。
% 全部题目按已建立的欧氏 Sobolev 工具自拟，不冒认未读教材的原题或题号。
% 固定非齐次能量 E_p(u)=int(|u|^p+|grad u|^p)，梯度取欧氏模，1<p<infty。
% 层次：1 为基础辨析；2--4 为容量比较与约束逼近，其中3为提高题；
% 5--6 为集合极限与反例；7--9 为局部化及细拓扑；10--11 为迹与边界空间。
% 不使用一般 Choquet 定理、容量版 Lebesgue 点定理或任意域的零拟处处刻画。

\begin{exercise}\label{b1:ex:26-basic-capacity-properties}
直接利用容量的定义和本节已经证明的结论，完成下列基本辨析。
\begin{enumerate}
 \item 若 \(A\subseteq B\subseteq\mathbb R^d\)，证明 \(C_p(A)\le C_p(B)\)。
 \item 对集合列 \((A_j)\)，写出容量的可列次可加不等式，并说明它不是可列可加等式。
 \item 选取一个固定的光滑截断函数，重新得到
 \[
  C_p(B(x,r))\le C_{d,p}(r^d+r^{d-p}).
 \]
 \item 当 \(1<p\le d\) 时，分别举出非空的零容量 Lebesgue 零测集与正容量 Lebesgue 零测集。
 当 \(p>d\) 时，证明不存在非空零容量集合。说明这些结论与 \(m_d^*(E)\le C_p(E)\) 的关系。
\end{enumerate}
\end{exercise}
% 来源：自拟基础辨析；只调用本节定义、次可加性、球估计和已证反例。
\begin{solution}
第一项由开邻域类的包含关系直接得到：每个包含 \(B\) 的开集也包含 \(A\)，先对容许函数、再对开邻域取下确界即可。

第二项是
\[
 C_p\left(\bigcup_{j=1}^{\infty}A_j\right)
 \le\sum_{j=1}^{\infty}C_p(A_j).
\]
左端即使对两两不交的集合也未必等于右端；本章后面的两点集合给出严格不等式。因此这里是可列次可加性，而不是测度的可列可加性。

取 \(\eta\in C_c^\infty(B(0,2))\)，使 \(0\le\eta\le1\) 且 \(\eta=1\) 于 \(B(0,1)\) 的一个邻域。函数
\[
 u_{x,r}(y)=\eta\left(\frac{y-x}{r}\right)
\]
在 \(B(x,r)\) 上容许，而变量代换给出
\[
 E_p(u_{x,r})
 =r^d\|\eta\|_p^p+r^{d-p}\|\nabla\eta\|_p^p.
\]
取只依赖 \(d,p,\eta\) 的常数便得第三项。

当 \(1<p\le d\) 时，由\cref{b1:prop:sobolev-point-capacity}，单点是非空、Lebesgue 零测且容量为零的集合。此时必有 \(d\ge2\)，平面片
\[
 [0,1]^{d-1}\times\{0\}
\]
的 \(d\) 维 Lebesgue 测度为零，但由\secref{b1:sec:2601}中竖线上的一维点值估计具有正容量。

当 \(p>d\) 时，对任意非空集合 \(E\) 取 \(x\in E\)。由单调性及\cref{b1:prop:sobolev-point-capacity}，
\[
 C_p(E)\ge C_p(\{x\})>0.
\]
因此不存在非空零容量集合；空集满足 \(C_p(\varnothing)=0\)。单点仍是零测而正容量的例子，特别包括 \(d=1\) 时本题允许的全部指数。不等式 \(m_d^*(E)\le C_p(E)\) 只说明零容量蕴含零测；它并未断言零测蕴含零容量，所以没有矛盾。
\end{solution}

\begin{exercise}\label{b1:ex:26-affine-capacity-comparison}
设 \(A\colon\mathbb R^d\rightarrow\mathbb R^d\) 为可逆线性映射，
\(b\in\mathbb R^d\)。证明对任意集合 \(E\subseteq\mathbb R^d\)，都有
\[
 |\det A|\min\{1,\|A\|^{-p}\}C_p(E)
 \le C_p(AE+b)
 \le|\det A|\max\{1,\|A^{-1}\|^p\}C_p(E).
\]
这里算子范数由欧氏范数诱导。由此证明容量在平移与正交变换下不变，并写出
\(E\mapsto tE\)、\(t>0\) 时的双边估计。解释为什么这些估计不同于齐次容量的单一幂次伸缩公式。
\end{exercise}
% 来源：自拟；把线性换元与非齐次能量结合，训练零阶项和梯度项的不同尺度。
\begin{solution}
先对光滑函数令 \(v(y)=u\bigl(A^{-1}(y-b)\bigr)\)。链式法则及换元给出
\[
 E_p(v)=|\det A|\int_{\mathbb R^d}
       \bigl(|u(x)|^p+|A^{-\mathsf T}\nabla u(x)|^p\bigr)\,\mathrm dx.
\]
对任意向量 \(\xi\)，有
\[
 \|A\|^{-1}\cdot|\xi|\le |A^{-\mathsf T}\xi|
                         \le\|A^{-1}\|\cdot|\xi|.
\]
因此所显示的两个常数分别给出 \(E_p(v)\) 相对于 \(E_p(u)\) 的下界和上界。
用全空间光滑稠密性逼近 \(u\)，这些估计同时保证复合函数及其候选弱梯度收敛，
所以同一公式和估计适用于所有 \(u\in W^{1,p}(\mathbb R^d)\)。

仿射映射把开集 \(G\) 双射到开集 \(AG+b\)，并保持 Lebesgue 零测集。
因而 \(u\ge1\) 几乎处处于 \(G\)，当且仅当相应的 \(v\ge1\) 几乎处处于 \(AG+b\)。
容许函数与开邻域在这两个映射下均一一对应。先对函数、再对开邻域取下确界，便得题设双边界；
不要求 \(E\) 可测。

当 \(A\) 正交时，两个常数都等于 \(1\)；取 \(A=I\) 同时得到平移不变性。
当 \(A=tI\) 时，估计化为
\[
 \min\{t^d,t^{d-p}\}C_p(E)
 \le C_p(tE)
 \le\max\{t^d,t^{d-p}\}C_p(E).
\]
零阶积分乘以 \(t^d\)，梯度积分乘以 \(t^{d-p}\)；一般不能把这两个因子合成一个相同的幂。
特别地，可逆仿射变换保持零容量集，但不保持一般集合的容量数值。
\end{solution}

\begin{optionalexercise}\label{b1:ex:26-one-dimensional-exact-capacity}
在 \(\mathbb R\) 上取 \(p=2\)，并设 \(L>0\)。计算
\[
 C_2(\{0\}),\quad C_2([0,L]),\quad C_2(\{0,L\}).
\]
证明两点集合的容量严格小于两个单点容量之和，由此说明 \(C_2\) 不是可加测度。
\end{optionalexercise}
% 来源：自拟；以一维能量配平方求精确值，区别非齐次容量与点质量的可加性。
\begin{solution}
先证半直线上的能量界。若 \(f\in H^1(0,\infty)\)，取其在有限区间上的绝对连续代表。
因为 \(f\in L^2(0,\infty)\)，可以选 \(R_j\to\infty\)，使 \(f(R_j)\to0\)。
在 \([0,R_j]\) 上展开平方并积分，有
\[
 \int_0^{R_j}(|f'|^2+|f|^2)\,\mathrm dx
 =\int_0^{R_j}|f'+f|^2\,\mathrm dx+|f(0)|^2-|f(R_j)|^2.
\]
令 \(j\to\infty\)，得到能量至少为 \(|f(0)|^2\)。反射后，左半直线也有相同结论。
在一维，\(H^1\) 函数的局部绝对连续代表连续；
若它在某点的开邻域几乎处处至少为 \(1\)，则该点值至少为 \(1\)。
因此，任意单点容许函数的总能量至少为 \(2\)。
函数 \(e^{-|x|}\) 的能量恰为 \(2\)，而
\((1+\varepsilon)e^{-|x|}\) 在原点的某个邻域大于 \(1\)。
让 \(\varepsilon\downarrow0\)，即得
\[
 C_2(\{0\})=2.
\]

对区间 \([0,L]\)，内部的零阶能量至少为 \(L\)，两侧的能量至少各为 \(1\)。
函数
\[
 v(x)=
 \begin{cases}
 e^x,&x<0,\\
 1,&0\le x\le L,\\
 e^{L-x},&x>L
 \end{cases}
\]
连续、分段光滑且属于 \(H^1(\mathbb R)\)，能量为 \(L+2\)。
用 \((1+\varepsilon)v\) 作开邻域容许函数，得到
\[
 C_2([0,L])=L+2.
\]

两点之间不再要求函数保持大于等于 \(1\)。令
\[
 h(x)=\frac{\cosh(x-L/2)}{\cosh(L/2)},\quad 0\le x\le L.
\]
它满足 \(h''=h\)、\(h(0)=h(L)=1\)，以及
\(h'(L)=-h'(0)=\tanh(L/2)\)。
任取两点集合的容许函数 \(u\)，在 \([0,L]\) 上写 \(u=h+w\)。
此时 \(w(0),w(L)\ge0\)，分部积分给出
\[
 \begin{aligned}
 \int_0^L(|u'|^2+|u|^2)\,\mathrm dx
 &=\int_0^L(|h'|^2+|h|^2)\,\mathrm dx
       +2[h'w]_0^L+\int_0^L(|w'|^2+|w|^2)\,\mathrm dx\\
 &\ge[h'h]_0^L=2\tanh(L/2).
 \end{aligned}
\]
加上左右两侧的能量下界，得到总能量至少为 \(2+2\tanh(L/2)\)。
将上面的 \(v\) 在中间区间改成 \(h\)，所得连续分段光滑函数达到这个能量；
再乘 \(1+\varepsilon\)，便在两个端点的开邻域满足容许条件。因此
\[
 C_2(\{0,L\})=2+2\tanh(L/2)<4=C_2(\{0\})+C_2(\{L\}).
\]
这里试探函数本身在端点取值恰为 \(1\)，乘以 \(1+\varepsilon\) 的步骤保证了容量定义要求的开邻域约束。
\end{solution}

\begin{exercise}\label{b1:ex:26-compact-smooth-admissibility}
设 \(K\subset\mathbb R^d\) 为紧集。证明
\[
 C_p(K)=\inf\left\{\,E_p(\varphi):
   \varphi\in C_c^\infty(\mathbb R^d),\
   0\le\varphi\le1,\
   \varphi=1\ \text{于 }K\text{ 的某个开邻域}\,\right\}.
\]
再在 \(1<p\le d\) 时，分别用一个无界可数集和一个有界稠密可数集，说明这个公式不能原样推广到任意集合。
\end{exercise}
% 来源：自拟；全空间稠密性的约束保持版本，新增远处截断和磨光尺度的选择。
\begin{solution}
记右端为 \(I(K)\)。其中每个光滑函数都是容量定义中的容许函数，所以 \(C_p(K)\le I(K)\)。
反向不等式不能仅由范数稠密性推出，因为任意光滑逼近未必保留邻域上的约束。

固定 \(\varepsilon>0\)。选开集 \(G\supseteq K\) 及容许函数 \(u\)，使
\[
 E_p(u)<C_p(K)+\varepsilon.
\]
这里紧集容量有限，可先用一个光滑截断函数得到有限上界。
将 \(u\) 截断到 \([0,1]\)，能量不增加，且 \(u=1\) 几乎处处于 \(G\)。
取 \(0\le\chi_R\le1\)，使 \(\chi_R=1\) 于 \(B(0,R)\)，
\(\operatorname{supp}\chi_R\subset B(0,2R)\)，并使
\(|\nabla\chi_R|\le C/R\)。乘积公式及 \(L^p\) 尾部收敛给出
\[
 \chi_R\cdot u\longrightarrow u\quad\text{于 }W^{1,p}(\mathbb R^d).
\]
确实，梯度差由 \((\chi_R-1)\cdot\nabla u\) 和 \(u\cdot\nabla\chi_R\) 两项组成；
前者由尾部积分趋零，后者的 \(L^p\) 范数至多为 \(C\|u\|_p/R\)。
选足够大的 \(R\)，使 \(K\subset B(0,R)\) 且
\(E_p(\chi_R\cdot u)<E_p(u)+\varepsilon\)。

紧集 \(K\) 到 \(\bigl(G\cap B(0,R)\bigr)^c\) 有正距离。
因此可选足够小的磨光尺度 \(\delta\)，使
\(\varphi=\rho_\delta*(\chi_R\cdot u)\) 在 \(K\) 的一个开邻域恒等于 \(1\)。
函数 \(\varphi\) 光滑紧支集，且 \(0\le\varphi\le1\)。
磨光对函数和向量梯度的 \(L^p\) 范数均为压缩，故
\[
 E_p(\varphi)\le E_p(\chi_R\cdot u)<C_p(K)+2\varepsilon.
\]
令 \(\varepsilon\downarrow0\)，证明了公式。

当 \(1<p\le d\) 时，可数集容量为零。对无界集合 \(\mathbb Z^d\)，
不存在紧支集函数能在其全体点的邻域恒等于 \(1\)，故对应的光滑下确界为无穷。
即使只考虑有界集合，也不能取消紧性条件。
取 \(E=\mathbb Q^d\cap B(0,1)\)。任何连续函数若在 \(E\) 的邻域恒等于 \(1\)，
则由稠密性在 \(\overline{B(0,1)}\) 上恒等于 \(1\)，能量至少为 \(m_d\bigl(B(0,1)\bigr)\)。
但 \(C_p(E)=0\)。后一个例子说明，连续容许函数会额外感受到集合的闭包。
\end{solution}

\begin{exercise}\label{b1:ex:26-capacity-hulls-and-decreasing-sets}
证明下列结论，并辨明最后一个极限为何不与第二项矛盾。
\begin{enumerate}
 \item\label{b1:item:26-capacity-gdelta-hull}
 每个集合 \(E\subseteq\mathbb R^d\) 都包含于一个 \(G_\delta\) 集 \(N\)，使
 \(C_p(N)=C_p(E)\)。特别地，零容量集有 Borel 零容量包络。
 \item\label{b1:item:26-capacity-decreasing-compact}
 若紧集列满足 \(K_1\supseteq K_2\supseteq\cdots\)，则
 \[
  C_p\left(\bigcap_{j=1}^\infty K_j\right)=\lim_{j\to\infty}C_p(K_j).
 \]
 \item\label{b1:item:26-capacity-decreasing-open}
 在 \(\mathbb R\) 上取 \(p=2\)。证明 \(G_j=(0,1/j)\) 递减到空集，而
 \(C_2(G_j)\to2\)。
\end{enumerate}
\end{exercise}
% 来源：自拟；从开外正则性推出限定的集合极限性质，不调用一般 Choquet 定理。
\begin{solution}
先设 \(C_p(E)<\infty\)。由外正则性选开集 \(O_j\supseteq E\)，使
\(C_p(O_j)<C_p(E)+1/j\)，并令 \(N=\bigcap_jO_j\)。
单调性给出
\[
 C_p(E)\le C_p(N)\le C_p(O_j)<C_p(E)+1/j.
\]
令 \(j\to\infty\)，得到第一项。若 \(C_p(E)=\infty\)，取 \(N=\mathbb R^d\) 即可。
这里仅断言存在同容量的 Borel 包络，没有断言容量在所有 Borel 集上可加。

令 \(K=\bigcap_jK_j\)。单调性给出 \(C_p(K)\le\lim_jC_p(K_j)\)。
任取包含 \(K\) 的开集 \(O\)。必有某个 \(j_0\)，使 \(K_j\subseteq O\) 对所有
\(j\ge j_0\) 成立：否则可以选 \(x_j\in K_j\setminus O\)，
由 \(K_1\) 紧性取收敛子列。固定任意 \(i\) 后，该子列的尾部落在闭集
\(K_i\setminus O\) 中，故其极限属于所有 \(K_i\) 却不属于 \(O\)，矛盾。
于是
\[
 \lim_jC_p(K_j)\le C_p(O).
\]
对开集 \(O\supseteq K\) 取下确界，即得第二项。

最后，\((0,1/j)\) 包含单点 \(1/(2j)\)，并包含于闭区间 \([0,1/j]\)。
由\cref{b1:ex:26-one-dimensional-exact-capacity}及平移不变性，
\[
 2\le C_2\bigl((0,1/j)\bigr)\le2+1/j.
\]
这些开区间的交为空集，容量却趋于 \(2\)。
它们不是紧集；其闭包都有共同点 \(0\)，而这个点未进入任何 \(G_j\)。
第二项证明中将极限留在每个闭集内的步骤，在这里恰好失效。
\end{solution}

\begin{exercise}\label{b1:ex:26-strong-convergence-without-full-pointwise-limit}
设 \(1<p<d\)。构造 \(u_n\in C_c^\infty(\mathbb R^d)\)，使
\[
 0\le u_n\le1,\quad u_n\longrightarrow0\ \text{于 }W^{1,p}(\mathbb R^d),
\]
但对每个 \(x\in[0,1]^d\)，都有
\[
 \liminf_{n\to\infty}u_n(x)=0,\quad
 \limsup_{n\to\infty}u_n(x)=1.
\]
由此说明范数收敛抽取拟处处收敛子列的结论，不能把“子列”删去。
\end{exercise}
% 来源：自拟移动峰；在每一层用有限个小立方体覆盖全部点，并插入零项固定下极限。
\begin{solution}
选取 \(0\le\eta\le1\) 的光滑函数，使它在 \([-1/2,1/2]^d\) 上恒等于 \(1\)，
支集包含于 \((-1,1)^d\)。对第 \(k\) 层的每个小立方体中心
\[
 c_{k,j}=2^{-k}\bigl(j+(1/2,\ldots,1/2)\bigr),
 \quad j\in\{0,\ldots,2^k-1\}^d,
\]
定义
\[
 v_{k,j}(x)=\eta\bigl(2^k(x-c_{k,j})\bigr).
\]
伸缩计算给出与 \(j\) 无关的界
\[
 E_p(v_{k,j})=
 2^{-kd}\|\eta\|_p^p+
 2^{-k(d-p)}\|\nabla\eta\|_p^p\longrightarrow0.
\]
按 \(k=1,2,\ldots\) 的次序列出每一层的全部 \(v_{k,j}\)，
并在相邻两项之间插入一个恒零函数，所得序列记为 \(u_n\)。
每层只有有限项，所以当 \(n\to\infty\) 时，非零项所在的层数趋于无穷。
上式因此证明了整列的 Sobolev 强收敛。

固定 \(x\in[0,1]^d\)。每一层都有一个闭小立方体包含 \(x\)，
它对应的函数在 \(x\) 处等于 \(1\)；立方体公共边界上的点也满足这一点。
故 \(u_n(x)=1\) 无穷多次，而插入的零项使 \(u_n(x)=0\) 也无穷多次。
结合 \(0\le u_n\le1\)，得到所要求的上下极限。
最后，由测度下界 \(C_p([0,1]^d)\ge1\)，这个不收敛集合不是零容量集。
所有 \(u_n\) 都是通常连续的代表，失败并非来自在零测集上任意改值。
\end{solution}

\begin{exercise}\label{b1:ex:26-local-quasicontinuous-gluing}
设 \(\Omega\subseteq\mathbb R^d\) 为开集。
对有限值函数 \(v\colon\Omega\rightarrow\mathbb R\)，把拟连续条件理解为：
每个 \(\varepsilon>0\) 都可找到全空间开集 \(G\)，满足 \(C_p(G)<\varepsilon\)，
且 \(v|_{\Omega\setminus G}\) 相对连续。
\begin{enumerate}
 \item\label{b1:item:26-local-quasicontinuous-cover}
 若 \(\Omega=\bigcup_jU_j\)，各 \(U_j\) 开，且 \(v|_{U_j}\) 均满足上述条件，
 则 \(v\) 在 \(\Omega\) 上拟连续。
 \item\label{b1:item:26-local-quasicontinuous-sobolev}
 每个 \(u\in W_{\mathrm{loc}}^{1,p}(\Omega)\) 都有有限值 Borel 拟连续代表，
 且两个这样的代表在 \(\Omega\) 上拟处处相等。
\end{enumerate}
\end{exercise}
% 来源：自拟局部化；使用光滑划分与全空间代表，不假定局部函数有有限的全空间能量。
\begin{solution}
对第一项，给定 \(\varepsilon>0\)，为第 \(j\) 个限制函数选坏开集 \(G_j\)，
使 \(C_p(G_j)<\varepsilon2^{-j-1}\)。集合 \(G=\bigcup_jG_j\) 的容量小于 \(\varepsilon\)。
固定 \(x\in\Omega\setminus G\)，选 \(j\) 使 \(x\in U_j\)。
在这个相对开邻域内，\(v\) 是连续函数
\(v|_{U_j\setminus G_j}\) 的限制，故 \(v\) 在 \(\Omega\setminus G\) 上连续。
连续性是局部性质，而统一的坏集由次可加性控制。

对第二项，取\cref{b1:thm:smooth-partition-unity}给出的非负光滑划分
\((\chi_j)\)，使其支集均紧含于 \(\Omega\)，在 \(\Omega\) 内局部有限，且
\(\sum_j\chi_j=1\)。
对每个 \(j\)，另取 \(\eta_j\in C_c^\infty(\Omega)\)，使
\(\eta_j=1\) 于 \(\operatorname{supp}\chi_j\) 的邻域。
函数 \(\eta_j\cdot u\) 的支集远离边界，由乘积公式和内部零延拓，
其全空间零延拓 \(u_j\) 属于 \(W^{1,p}(\mathbb R^d)\)。
为它选一个有限值 Borel 拟连续代表 \(\widetilde u_j\)，并在 \(\Omega\) 上令
\[
 v(x)=\sum_j\chi_j(x)\cdot\widetilde u_j(x).
\]
局部有限性保证这个和在每点附近都是有限和，因此有限且 Borel 可测。
在共同的满测度集合上，各 \(\widetilde u_j=u_j\)，于是
\[
 v=\sum_j\chi_j\cdot\eta_j\cdot u=\sum_j\chi_j\cdot u=u
 \quad\text{几乎处处于 }\Omega.
\]
再为各 \(\widetilde u_j\) 选容量小于 \(\varepsilon\cdot2^{-j-1}\) 的坏开集。
它们的并以外，每项 \(\chi_j\cdot\widetilde u_j\) 连续，且和仍局部有限，
故 \(v\) 拟连续。

设 \(v,w\) 是 \(u\) 的两个有限值拟连续代表。
对每个 \(j\)，将 \(\chi_j\cdot(v-w)\) 在 \(\Omega\) 外补零。
这个函数全空间拟连续：在所选坏开集外，乘积在 \(\Omega\) 内连续；
而 \(\operatorname{supp}\chi_j\) 紧含于 \(\Omega\)，所以在每个边界点附近乘积恒为零。
它几乎处处等于零，是零 Sobolev 类的拟连续代表。
由\cref{b1:thm:quasicontinuous-uniqueness}，存在容量零集 \(N_j\)，使
\(\chi_j\cdot(v-w)=0\) 于 \(N_j^c\)。
删去 \(N=\bigcup_jN_j\) 后，对每点至少有一个 \(\chi_j>0\)，从而 \(v=w\)。
证明没有把仅具局部能量的 \(u\) 直接代入全空间超水平估计。
\end{solution}

\begin{optionalexercise}\label{b1:ex:26-zero-capacity-thinness}
设 \(N\subseteq\mathbb R^d\) 且 \(C_p(N)=0\)。
\begin{enumerate}
 \item 证明 \(N\) 在每一点都是 \(p\)-薄的，并由此证明 \(\mathbb R^d\setminus N\) 是细开集。
 \item 设 \(d\ge2\)、\(1<p\le d\)，并令
 \(\Omega=B(0,1)\setminus\{0\}\)。直接判断补集在 \(0\) 处的 Wiener 级数是收敛还是发散；再利用选读的 Wiener 判据判断 \(0\) 是否为正则边界点。
\end{enumerate}
\end{optionalexercise}
% 来源：自拟；训练全空间零容量到相对薄性的截断接口及穿孔球判据。
\begin{solution}
固定 \(x\) 与 \(r>0\)。取一个在 \(B(x,r)\) 上等于一、支集含于 \(B(x,2r)\) 的光滑截断 \(\eta\)。若 \(u\) 是覆盖 \(N\) 的全空间容许函数，则 \(\eta\cdot u\) 在 \(B(x,2r)\) 中零边界，并覆盖 \(N\cap B(x,r)\)。乘积估计给出
\[
 \operatorname{cap}_p\bigl(N\cap B(x,r),B(x,2r)\bigr)
 \le C_{x,r,p}\cdot E_p(u).
\]
依次对容许函数和开邻域取下确界，右端为零。故薄性定义中的每一项都为零，\(N\) 在每一点都薄。对 \(N\) 的补集应用细开集定义，即得第一项。

对穿孔球，在所有充分小的半径上，
\[
 (\mathbb R^d\setminus\Omega)\cap B(0,r)=\{0\}.
\]
单点全空间容量为零，第一项的论证说明相应相对容量也为零，故 Wiener 级数收敛，补集在 \(0\) 处是薄的。选读的 Wiener 判据因此说明 \(0\) 不是正则边界点。
\end{solution}

\begin{optionalexercise}\label{b1:ex:26-exterior-corkscrew-wiener}
设 \(d\ge2\)、\(1<p\le d\)，\(\Omega\subset\mathbb R^d\) 为有界开集，且 \(x_0\in\partial\Omega\)。假设存在 \(a\in(0,1/4)\) 与 \(r_0>0\)，使每个 \(0<r<r_0\) 都有
\[
 B(y_r,ar)\subseteq B(x_0,r)\setminus\Omega.
\]
用零边界 Poincaré 不等式证明
\[
 \operatorname{cap}_p\bigl(B(x_0,r)\setminus\Omega,B(x_0,2r)\bigr)
 \ge c_{d,p,a}r^{d-p}.
\]
由此验证 Wiener 级数发散；再利用选读的 Wiener 判据说明 \(x_0\) 正则。
\end{optionalexercise}
% 来源：自拟；把外部软木塞条件直接转成 Wiener 级数的尺度下界。
\begin{solution}
设 \(v\in W_0^{1,p}(B(x_0,2r))\) 在包含 \(B(y_r,ar)\) 的开集上几乎处处至少为一。截断后可设 \(0\le v\le1\)，于是
\[
 \omega_d a^dr^d
 \le\int_{B(x_0,2r)}|v|^p
 \le (4r)^p\int_{B(x_0,2r)}|\nabla v|^p,
\]
其中第二个不等式是零边界 Poincaré 不等式。对所有容许函数取下确界，并用内集单调性，得到题设下界。

当 \(r=2^{-j}<r_0\) 时，Wiener 级数的第 \(j\) 项至少为正常数
\(c_{d,p,a}^{1/(p-1)}\)。因此级数发散，补集在 \(x_0\) 处不是薄的；选读的 Wiener 判据给出 \(x_0\) 正则。
\end{solution}

\begin{exercise}\label{b1:ex:26-quasicontinuous-plane-trace}
设 \(d\ge2\)，\(P=\mathbb R^{d-1}\times\{0\}\)，以 \(\sigma\) 表示 \(P\) 上的
\((d-1)\) 维 Lebesgue 测度，\(\sigma^*\) 表示相应外测度。
\begin{enumerate}
 \item\label{b1:item:26-plane-capacity-control}
 证明存在仅依赖 \(p\) 的常数 \(A_p\)，使对任意集合 \(E\subseteq\mathbb R^d\)，
 \[
  \sigma^*(E\cap P)\le A_p\cdot C_p(E).
 \]
 \item\label{b1:item:26-plane-trace-representative}
 设 \(u\in W^{1,p}(\mathbb R^d)\)，\(\widetilde u\) 是任一有限值拟连续代表。
 证明 \(\widetilde u|_P\) 与上半空间限制的 \(L^p\) 迹
 \(\operatorname{Tr}_P u\) 在 \(\sigma\)-几乎处处意义下相等。
 因而迹虽然不能由任意 Lebesgue 代表的平面点值给出，却可由拟连续代表给出。
\end{enumerate}
\end{exercise}
% 来源：自拟接口题；ACL 的一维点值界控制平面测度，再以拟处处逼近辨认迹。
\begin{solution}
对绝对连续函数 \(f\colon[0,1]\rightarrow\mathbb R\)，先由基本定理对终点平均，得到
\[
 |f(0)|\le\int_0^1|f(t)|\,\mathrm dt+\int_0^1|f'(t)|\,\mathrm dt.
\]
Hölder 不等式及标量幂次估计给出
\[
 |f(0)|^p\le2^{p-1}\int_0^1\bigl(|f(t)|^p+|f'(t)|^p\bigr)\,\mathrm dt.
\]
沿几乎每条竖线使用 ACL 代表，再对 \(x'\in\mathbb R^{d-1}\) 积分，
便得到平面迹的界
\[
 \|\operatorname{Tr}_P u\|_{L^p(P)}^p\le A_p\cdot E_p(u),
 \quad A_p=2^{p-1}.
\]
这个迹就是\cref{b1:thm:half-space-lp-trace}的端点迹；
该不等式也说明它对全空间 Sobolev 范数连续。

设 \(G\supseteq E\) 开，\(u\) 是 \(G\) 的容量容许函数。
Fubini 定理允许先排除一个横向零测集，使 \(u\ge1\) 在每个相关竖截面内几乎处处成立，
并使这些竖线上的 ACL 代表都存在。
若 \((x',0)\in G\)，开性保证该竖截面含有原点的一个区间。
沿线连续性于是推出 \(\operatorname{Tr}_P u(x')\ge1\)。
所以
\[
 \sigma^*(E\cap P)\le\sigma(G\cap P)
 \le\|\operatorname{Tr}_P u\|_{L^p(P)}^p
 \le A_p\cdot E_p(u).
\]
依次对容许函数及开邻域取下确界，得到第一项。

由全空间光滑稠密性及\cref{b1:thm:sobolev-quasi-everywhere-subsequence}，
可选 \(\varphi_j\in C_c^\infty(\mathbb R^d)\)，使
\(\varphi_j\to u\) 于 \(W^{1,p}\)，并拟处处收敛到 \(\widetilde u\)。
第一项说明这个收敛的例外集在 \(P\) 上具有 \(\sigma\) 外测度零。
另一方面，迹的连续性给出
\[
 \varphi_j|_P=\operatorname{Tr}_P\varphi_j
       \longrightarrow\operatorname{Tr}_P u
       \quad\text{于 }L^p(P).
\]
再抽一个子列，使平面上的收敛几乎处处成立。
同一子列的两个点态极限必须相同，所以
\(\widetilde u|_P=\operatorname{Tr}_P u\) 在 \(\sigma\)-几乎处处成立。
这也保证了限制函数在平面完备化测度下可测。
若只指定任意 Lebesgue 代表，则可以在整个零体积的平面上随意改值，
因而不能得到同样结论。
\end{solution}

\begin{exercise}\label{b1:ex:26-slit-zero-boundary-space}
设 \(d\ge2\)，
\[
 \Omega=(-2,2)^d,\quad
 K=[-1,1]^{d-1}\times\{0\},\quad D=\Omega\setminus K.
\]
证明下列结论。
\begin{enumerate}
 \item\label{b1:item:26-slit-positive-capacity}
 集合 \(K\) 的 \(d\) 维 Lebesgue 测度为零，但 \(C_p(K)>0\)。
 \item\label{b1:item:26-slit-strict-w0}
 经零延拓并限制到 \(\Omega\)，空间 \(W_0^{1,p}(D)\) 是
 \(W_0^{1,p}(\Omega)\) 的闭真子空间。
 \item\label{b1:item:26-slit-zero-extension-insufficient}
 存在 \(v\in W^{1,p}(D)\)，使其全空间零延拓属于
 \(W^{1,p}(\mathbb R^d)\)，但 \(v\notin W_0^{1,p}(D)\)。
\end{enumerate}
解答中只用 \(W_0^{1,p}\) 的闭包定义及已经建立的迹，不调用一般域上的零拟处处刻画。
\end{exercise}
% 来源：自拟裂缝域；零体积障碍仍可留下非零迹约束，连接容量与闭包定义。
\begin{solution}
集合 \(K\) 位于一个超平面内，故 \(m_d(K)=0\)，而
\(\sigma(K)=2^{d-1}\)。
由\cref{b1:ex:26-quasicontinuous-plane-trace}，
\[
 C_p(K)\ge A_p^{-1}2^{d-1}>0.
\]

若 \(f\in W_0^{1,p}(D)\)，取 \(f_j\in C_c^\infty(D)\) 在该空间范数中逼近 \(f\)。
每个 \(f_j\) 在裂缝及 \(\partial\Omega\) 附近均为零，因此在 \(\Omega\) 上补零后属于
\(C_c^\infty(\Omega)\)。
由零延拓\cref{b1:prop:w0-zero-extension}，这些延拓在全空间 Sobolev 范数中收敛，
其极限限制到 \(\Omega\) 属于 \(W_0^{1,p}(\Omega)\)。
由于 \(K\) 的体积为零，这个映射保持 Sobolev 范数。
源空间完备，故它的像是闭子空间。

为了证明像不是整个空间，选 \(\varphi\in C_c^\infty(\Omega)\)，使
\(\varphi=1\) 于 \(K\) 的某个邻域，并令 \(v=\varphi|_D\)。
假设 \(v\in W_0^{1,p}(D)\)，则存在 \(f_j\in C_c^\infty(D)\) 满足
\(f_j\to v\) 于 \(W^{1,p}(D)\)。
将 \(f_j\) 延拓为全空间光滑紧支集函数；由于 \(m_d(K)=0\)，有
\[
 \|f_j-\varphi\|_{W^{1,p}(\mathbb R^d)}
       =\|f_j-v\|_{W^{1,p}(D)}\longrightarrow0.
\]
平面迹连续，故 \(\operatorname{Tr}_P f_j\to\operatorname{Tr}_P\varphi\) 于 \(L^p(P)\)。
但各 \(f_j\) 在 \(K\) 上恒为零，而 \(\varphi|_K=1\)，于是
\[
 \|\operatorname{Tr}_P f_j-\operatorname{Tr}_P\varphi\|_{L^p(P)}^p
       \ge \sigma(K)=2^{d-1},
\]
矛盾。故 \(\varphi\) 不在上述像中，第二项得证。

同一个 \(v\) 也证明第三项：它在 \(D\) 外补零所得函数与全空间的
\(\varphi\) 仅在 \(K\) 上可能不同，因而是同一个 \(W^{1,p}(\mathbb R^d)\) 类；
但是刚才已经证明 \(v\notin W_0^{1,p}(D)\)。
零延拓属于 Sobolev 空间，只约束体积意义下的弱导数；
裂缝上的额外零迹约束则由内部光滑紧支集逼近保留下来。
\end{solution}
